% !TeX root = ../../Book1.tex
% 纲要标题：## 第1章　仿射、Euclid 结构与变换群
% 纲要位置：几何学计划纲要.md，第 117 行。
% 正文状态：初稿；尚未正式审核。

% 纲要内容开始
% 
% > 本章目标与证明范围：以“选择什么结构、允许什么变换、分类保留什么不变量”为统一问题。完整处理仿射构造、实二次型和 Euclid 等距变换；群作用在实际分类中建立，不另作一章名词准备。
% 
% 纲要内容结束

\chapter{仿射、Euclid 结构与变换群}
\label{g1:ch:01}
\BookChapterNumber{g1:ch:01}

\begin{chapterabstract}
	几何分类取决于指定的结构与允许的变换。本章从点与位移的区别建立仿射空间，说明坐标、重心与子空间相交的内在含义；加入内积后，证明正交投影、自伴谱分解、实对称型的惯性分类，并用内积矩阵的行列式计算体积。随后从距离保持推出 Euclid 等距变换的仿射形式，分解其正交部分与不可消去的平移，证明反射生成并给出低维分类。最后以轨道和稳定子群统一这些分类，比较 Euclid 、相似与仿射等价，考察退化参数处不变量的变化。
\end{chapterabstract}

\begin{learninggoals}
	\begin{itemize}
		\item 区分点与向量，计算仿射组合、标架坐标及仿射子空间的相交与包络。
		\item 说明正交投影和谱分解的构造，区分正交相似、一般相似与合同，求实对称型的惯性，并由内积矩阵计算体积。
		\item 从距离条件求出等距变换的正交部分和平移部分，判断不动点、最小位移及低维运动类型。
		\item 用轨道、稳定子群和完整分类参数比较不同等价关系，并判断连续变化何时跨越退化对象。
	\end{itemize}
\end{learninggoals}

把一个三角形在桌面上平移或旋转，边长没有改变；把它画到另一个比例的地图上，边长改变而比例没有改变；把图形沿一个方向拉长，连角度也可能改变。但只要三个顶点始终不共线，它们仍然确定同一种仿射标架，线段的中点也仍被送到中点。这三个操作分别提醒我们：长度、长度比和仿射组合是不同层次的信息。说两个图形“相同”之前，必须说明我们希望保存其中哪些信息。

\ProseParagraphBreak

本章将这一问题落实到可计算的构造中。先不选长度，只研究点之间的位移；再在方向空间上指定内积，使长度、正交和体积获得意义。二次型提供了观察“同一数据在不同坐标中怎样表达”的具体机会，而 Euclid 运动则让我们看到，仅由距离也能恢复仿射结构。旋转中心、螺旋轴或滑移方向会从同一个线性方程中出现，群作用随后把这些计算概括为分类问题。

\ProseParagraphBreak

全文限于有限维实空间。所需基础是向量空间、基、线性映射、矩阵、行列式与秩的基本运算；关键的几何用法和谱定理本身在本章证明。另用有限维空间中闭有界集的紧性，以及连续实函数在非空紧集上取得最大、最小值这一基本分析结论。例如在单位球面上取函数值趋近上确界的点列，紧性给出收敛子列，连续性便使极限点取得该上确界；这一短论证正是后文极值法的依据。

\ProseParagraphBreak

仿射构造可对照 Gallier 的正式第二版\cite[第2.1--2.8, 2.10--2.11节]{Gallier2011Geometric}；内积、等距变换与矩阵标准形可分别参看同书\cite[第6.1--6.4, 8.1--8.2, 8.4--8.7, 12.3--12.4节]{Gallier2011Geometric}。

% !TeX root = ../../Book1.tex
% 纲要标题：### 1.1 仿射空间的内在构造
% 纲要位置：几何学计划纲要.md，第 121 行。
% 正文状态：初稿；尚未正式审核。

% 纲要内容开始
% 
% * 平移作用、点差与仿射组合
% * 仿射标架、坐标变换与重心坐标
% * 子空间的相交、平行与维数公式
% 
% 纲要内容结束

\section{仿射空间的内在构造}
\label{g1:sec:0101}

在一张地图上改变坐标原点，不会改变两个地点之间的位移，却会同时改变两个地点的坐标。因此，“地点”与“位移”应当是两类对象。向量空间天然带有零向量，仿射空间则把原点的选择从几何结构中分离出来。这样做以后，中点、直线和平行关系仍可定义，而两个点直接相加通常没有意义。

\subsection{平移作用、点差与仿射组合}
\label{g1:subsec:0101:01}

\begin{definition}[仿射空间]\label{g1:def:affine-space}
	设 \(V\) 是有限维实向量空间。给非空集合 \(A\) 指定运算 \(A\times V\rightarrow A\)，记为 \((p,v)\mapsto p+v\)。若该运算满足
	\[
		p+0=p,\qquad (p+u)+v=p+(u+v),
	\]
	且对每对 \(p,q\in A\)，恰有一个 \(v\in V\) 满足 \(p+v=q\)，则称这个带有运算的集合为以 \(V\) 为方向空间的\term[fangshekongjian]{仿射空间}[affine space]。满足上述等式的唯一向量称为\term[diancha]{点差}[difference vector]，记为 \(q-p\)；定义 \(\dim A=\dim V\)。
\end{definition}

映射 \(\tau_v:p\mapsto p+v\) 称为\term[pingyi]{平移}[translation]。上述条件说，每一个点都能通过唯一的平移向量移到另一个点。特别地，\(q=p+(q-p)\)，而点差满足
\[
	(r-q)+(q-p)=r-p,\qquad p-q=-(q-p).
\]
例如，第一个等式由 \(p+((q-p)+(r-q))=q+(r-q)=r\) 及位移的唯一性得到。这里的加法发生在方向空间中，不能把它误读成三个点的普通加减。

\begin{example}[方程的解集]\label{g1:exm:affine-solutions}
	设 \(L:V\rightarrow W\) 为实线性映射，方程 \(Lx=b\) 有一个解 \(p\)。其解集恰为 \(p+\ker L\)：若 \(x\) 也是解，则 \(L(x-p)=0\)；若 \(u\in\ker L\)，则 \(L(p+u)=b\)。它是以 \(\ker L\) 为方向空间的仿射空间。若 \(b\ne0\)，这个解集不含零向量，不能当作 \(V\) 的向量子空间。例如 \(x+y+z=1\) 是仿射平面，两个解之和一般已不是解。
\end{example}

选定 \(o\in A\) 后，\(p\mapsto p-o\) 将 \(A\) 与 \(V\) 一一对应。若改选 \(o'=o+a\)，则同一点的向量表示变为 \(p-o'=(p-o)-a\)。这说明可以借助原点计算，但要检查计算结果是否依赖它。特别值得保留的运算是系数和为 \(1\) 的加权平均。

\begin{proposition}[仿射组合]\label{g1:prop:affine-combination}
	设 \(A\) 是实仿射空间，\(p_0,\ldots,p_m\in A\)， \(\lambda_0,\ldots,\lambda_m\in\mathbb R\) 且 \(\sum_i\lambda_i=1\)。则
	\[
		o+\sum_{i=0}^m\lambda_i(p_i-o)
	\]
	与 \(o\in A\) 的选择无关，称为这些点的\term[fangshezuhe]{仿射组合}[affine combination]，简记为 \(\sum_i\lambda_i p_i\)。
\end{proposition}
\begin{proof}
	以 \(o'=o+a\) 代替 \(o\)，新表达式为
	\[
		o+a+\sum_i\lambda_i\bigl((p_i-o)-a\bigr)
		=o+\sum_i\lambda_i(p_i-o)+(1-\sum_i\lambda_i)a.
	\]
	最后一项为零，因而得到同一点。
\end{proof}

系数和条件正是消去原点影响的条件。特别地，当 \(p\ne q\) 时，\((1-t)p+tq=p+t(q-p)\) 描述穿过 \(p,q\) 的直线；当 \(0\le t\le1\) 时描述线段，当 \(t=1/2\) 时给出中点。若 \(p=q\)，这些组合都退化为同一点。若系数和为零，则 \(\sum_i\lambda_i(p_i-o)\) 也与原点无关，但它给出的是向量而不是点。这一差别解释了为什么 \(q-p\) 合法，而 \(p+q\) 不能作为不依赖原点的点运算。

\begin{definition}[仿射映射]\label{g1:def:affine-map}
	设 \(A,B\) 分别以实向量空间 \(V,W\) 为方向空间。对于映射 \(f:A\rightarrow B\)，若存在实线性映射 \(L:V\rightarrow W\)，使
	\[
		f(p+v)=f(p)+Lv\qquad(p\in A,\ v\in V),
	\]
	则称这个映射为\term[fangsheyingshe]{仿射映射}[affine map]，称 \(L\) 为 \(f\) 的\term[xianxingbufen]{线性部分}[linear part]。
\end{definition}

线性部分唯一，因为 \(Lv=f(p+v)-f(p)\)。选定原点后，仿射映射就是 \(x\mapsto Lx+b\)。如果 \(L\) 可逆，逆映射是 \(y\mapsto L^{-1}(y-b)\)，此时称 \(f\) 为\term[fangshetonggou]{仿射同构}[affine isomorphism]；反过来，若仿射映射是双射，则其线性部分也是双射。

\begin{proposition}\label{g1:prop:affine-preservation}
	实仿射空间间的映射 \(f:A\rightarrow B\) 是仿射映射，当且仅当它保持一切有限仿射组合。
\end{proposition}
\begin{proof}
	若 \(f\) 的线性部分为 \(L\)，在组合的计算中以 \(o\) 与 \(f(o)\) 为原点，便有
	\[
		f\left(o+\sum_i\lambda_i(p_i-o)\right)
		=f(o)+\sum_i\lambda_i\bigl(f(p_i)-f(o)\bigr).
	\]
	这正是保持仿射组合的等式。

	反过来，固定 \(o\)，定义 \(L(v)=f(o+v)-f(o)\)。从组合 \((1-t)o+t(o+v)\) 得到 \(L(tv)=tL(v)\)，其中 \(t\) 为任意实数。再对三点 \(o+u,o+v,o\) 使用系数 \(1,1,-1\)，得到 \(L(u+v)=L(u)+L(v)\)。故 \(L\) 线性，且 \(f(p+v)=f(o)+L(p-o+v)=f(p)+Lv\)。
\end{proof}

\subsection{仿射标架、坐标变换与重心坐标}
\label{g1:subsec:0101:02}

要给一个点写出数值坐标，需要选原点与方向基；要用若干几何点描述它，则可以用一个有序点组同时指定原点与方向基。

\begin{definition}[标架与仿射无关]\label{g1:def:affine-frame}
	在 \(n\) 维实仿射空间 \(A\) 中，若 \(p_1-p_0,\ldots,p_m-p_0\) 线性无关，就称 \(p_0,\ldots,p_m\) \term[fangshewuguan]{仿射无关}[affinely independent]。当 \(m=n\) 时，有序点组 \((p_0,\ldots,p_n)\) 称为\term[fangshebiaojia]{仿射标架}[affine frame]。
\end{definition}

仿射无关不依赖于把哪个点放在首位。事实上，它等价于：若 \(\sum_i c_i=0\) 且 \(\sum_i c_i(p_i-o)=0\)，则所有 \(c_i\) 为零。取 \(o=p_0\) 就得到原定义；这个对称表述不再区别各点。

\begin{proposition}[坐标及其变换]\label{g1:prop:affine-coordinates}
	设 \(A\) 为 \(n\) 维实仿射空间。给定仿射标架 \((p_0,\ldots,p_n)\)，每个 \(p\in A\) 唯一写成
	\[
		p=p_0+\sum_{i=1}^n x_i(p_i-p_0)
		=\sum_{i=0}^n\lambda_i p_i,\qquad
		\lambda_0=1-\sum_{i=1}^n x_i,\quad\lambda_i=x_i\ (i\ge1).
	\]
	称 \((x_1,\ldots,x_n)\) 为\term[fangshezuobiao]{仿射坐标}[affine coordinates]，\((\lambda_0,\ldots,\lambda_n)\) 为相对于此标架的\term[zhongxinzuobiao]{重心坐标}[barycentric coordinates]。若新原点相对于旧标架的坐标为 \(a\)，新方向基在旧方向基中的列矩阵为 \(S\)，则新旧坐标满足
	\[
		x=a+Sx',\qquad x'=S^{-1}(x-a).
	\]
\end{proposition}
\begin{proof}
	基展开给出 \(p-p_0\) 的唯一系数 \(x_i\)，而系数和为 \(1\) 唯一确定 \(\lambda_0\)。设旧方向基的组合记为 \(Ex=\sum_i x_i(p_i-p_0)\)，则新原点为 \(p_0+Ea\)，新方向基对应线性映射 \(ES\)。于是 \(p=p_0+Ex=p_0+Ea+ESx'\)。由于 \(E\) 单射，得到 \(x=a+Sx'\)；新方向基仍是基，所以 \(S\) 可逆。
\end{proof}

\begin{example}[三角形中的位置]\label{g1:exm:barycentric-triangle}
	取 \(p_0=(0,0),p_1=(3,0),p_2=(0,2)\)。点 \(p=(1,1)\) 的重心坐标为 \((1/6,1/3,1/2)\)，因为后两个系数分别由 \(3\lambda_1=1\)、 \(2\lambda_2=1\) 决定，再由系数和求出第一个。三个系数全非负表示点处在闭三角形内，全为正表示点处在内部；在全非负的条件下，某个系数为零表示点处在对应的边上。例如 \(q=(2,1)\) 的坐标为 \((-1/6,2/3,1/2)\)，它越过了与 \(p_0\) 相对的边。
\end{example}

\begin{figure}[htbp]
	\centering
	% !TeX root = ../Book1.tex
\begin{tikzpicture}[x=1.45cm,y=1.45cm]
	\fill[GeometryBlueLight] (0,0)--(3,0)--(0,2)--cycle;
	\draw[thick] (0,0)--(3,0)--(0,2)--cycle;
	\draw[dashed,gray] (-.3,2.2)--(3.3,-.2);
	\fill (0,0) circle (1.5pt) node[below left] {\(p_0\)};
	\fill (3,0) circle (1.5pt) node[below right] {\(p_1\)};
	\fill (0,2) circle (1.5pt) node[above left] {\(p_2\)};
	\fill[GeometryBlue] (1,1) circle (1.7pt) node[below left] {\(p=(1,1)\)};
	\fill[GeometryBlue] (2,1) circle (1.7pt) node[right] {\(q=(2,1)\)};
	\node at (1,.24) {\(\lambda_0>0\)};
\end{tikzpicture}

	\caption[重心坐标与三角形]{重心坐标与三角形。直线 \(p_1p_2\) 上 \(\lambda_0=0\)，点 \(p\) 与顶点 \(p_0\) 在同侧，点 \(q\) 在另一侧。}
	\label{g1:fig:barycentric}
\end{figure}

图中的符号判断可以直接计算，而无须凭图形猜测：此例 \(\lambda_0=1-x/3-y/2\)，所以直线两侧的符号由一次函数决定。一般情形下，系数全非负的仿射组合构成这组仿射无关点张成的\term[dan-chun-xing]{单纯形}[simplex]。这种位置关系在任何仿射坐标中都成立，因为仿射同构及其逆映射都保持这些组合，而不要求保持长度或角度。

\begin{proposition}[由标架确定映射]\label{g1:prop:frame-map}
	设 \((p_0,\ldots,p_n)\) 为实仿射空间 \(A\) 的标架，\(B\) 为另一个实仿射空间。任给 \(q_0,\ldots,q_n\in B\)，存在唯一仿射映射 \(f:A\rightarrow B\) 满足 \(f(p_i)=q_i\)。若 \(B\) 也为 \(n\) 维，则 \(f\) 是同构当且仅当这些 \(q_i\) 构成标架。
\end{proposition}
\begin{proof}
	定义线性映射 \(L\) 将基向量 \(p_i-p_0\) 送到 \(q_i-q_0\)，再置 \(f(p)=q_0+L(p-p_0)\)。它是满足指定条件的仿射映射。任何满足条件的映射都必须有这一线性部分及同一 \(f(p_0)\)，故唯一。两空间同维时，\(L\) 可逆恰好要求其基向量的像构成基。
\end{proof}

\subsection{子空间的相交、平行与维数公式}
\label{g1:subsec:0101:03}

在向量空间中，子空间一定交于零向量；在仿射空间中，两条直线可能平行，也可能既不相交又不平行。为了区别这些情形，除了方向空间，还要记录两个空间之间的位移。

\begin{definition}[仿射子空间]\label{g1:def:affine-subspace}
	设 \(A\) 为以 \(V\) 为方向空间的实仿射空间，\(p\in A\)，\(U\subseteq V\) 为向量子空间。集合 \(F=p+U\) 称为\term[fangshezikongjian]{仿射子空间}[affine subspace]，其方向空间为 \(U\)，维数为 \(\dim U\)。若两个仿射子空间的方向空间之一包含于另一个，本章就称它们\term[pingxing]{平行}[parallel]；对同维子空间，这等价于两个方向空间相等。允许重合的空间互相平行。
\end{definition}

不同文献对异维空间的“平行”用语不完全一致，后文的判断将直接写出方向空间的包含或相等。对于 \(F=p+U\)，选其中任何 \(r\) 为基点仍有 \(F=r+U\)，因为 \(r-p\in U\)。同时 \(F-F=U\)，所以方向空间由点集本身唯一决定。

\begin{proposition}[相交与仿射包络]\label{g1:prop:affine-intersection}
	在有限维实仿射空间 \(A\) 中，设 \(F=p+U\)、\(G=q+W\) 为非空仿射子空间。则
	\[
		F\cap G\ne\varnothing\quad\Longleftrightarrow\quad q-p\in U+W.
	\]
	若 \(r\in F\cap G\)，则 \(F\cap G=r+(U\cap W)\)。包含 \(F\cup G\) 的最小仿射子空间，称为其\term[fangshebaoluo]{仿射包络}[affine hull]，为
	\begin{equation}\label{g1:eq:affine-join}
		\operatorname{aff}(F\cup G)=p+\bigl(U+W+\mathbb R(q-p)\bigr).
	\end{equation}
	特别地，令 \(d=\dim(U\cap W)\)，则
	\[
		\dim\operatorname{aff}(F\cup G)
		=\begin{cases}
			\dim F+\dim G-d,&F\cap G\ne\varnothing,\\
			\dim F+\dim G-d+1,&F\cap G=\varnothing.
		\end{cases}
	\]
	在第一种情形中，\(d=\dim(F\cap G)\)。
\end{proposition}
\begin{proof}
	交点的方程为 \(p+u=q+w\)，也就是 \(q-p=u-w\)。这个方程有解，当且仅当 \(q-p\in U+W\)。若有交点 \(r\)，则写成 \(F=r+U\)、\(G=r+W\)，一个点属于二者等价于它与 \(r\) 的差同时属于 \(U,W\)，给出交集公式。

	式\eqref{g1:eq:affine-join}右边包含 \(F,G\)。任何包含二者的仿射子空间都含 \(p\)，其方向空间必须包含 \(U,W\) 及 \(q-p\)，因而也包含右边的方向空间。这证明了最小性。

	最后回顾线性空间的维数计算。先取 \(U\cap W\) 的一组基，分别扩充成 \(U,W\) 的基；把交集基只写一次，再合并两个扩充部分，就得到 \(U+W\) 的基。张成性来自两个扩充基；若合并后有线性关系，则 \(U\) 的扩充部分之和落在 \(W\)，因而落在交集，这迫使该部分系数为零，再依次得到其余系数为零。所以 \(\dim(U+W)=\dim U+\dim W-d\)。交点存在时 \(q-p\) 已在这个和中；交点不存在时再加入它恰增加一维。
\end{proof}

\begin{example}[平行线与异面直线]\label{g1:exm:skew-lines}
	在 \(\mathbb R^3\) 中，令
	\[
		F=\{(s,0,0):s\in\mathbb R\},\qquad
		G=\{(0,t,1):t\in\mathbb R\}.
	\]
	它们的方向分别为 \(\mathbb Re_1,\mathbb Re_2\)，基点之差为 \(e_3\)，不在 \(\mathbb Re_1+\mathbb Re_2\) 中，故不相交；方向也不平行。两条这样的直线称为\term[yimianzhixian]{异面直线}[skew lines]。其仿射包络维数为 \(1+1-0+1=3\)。若改成 \(G'=\{(t,0,1):t\in\mathbb R\}\)，方向交集为一维，包络维数变为 \(2\)，这才是两条不同的平行线。
\end{example}

对任意非空点集 \(S\subseteq A\)，选定 \(p_0\in S\) 后，同一最小性论证给出
\[
	\operatorname{aff}(S)=p_0+\operatorname{span}\{p-p_0:p\in S\}.
\]
这里的 \(\operatorname{span}\) 只取有限线性组合。展开右边后可见，它也恰好是 \(S\) 中有限多个点的一切仿射组合。因此有限点组 \(p_0,\ldots,p_m\) 的仿射包络维数就是列向量 \(p_1-p_0,\ldots,p_m-p_0\) 的秩，是否仿射无关可以转化为一次矩阵秩计算。

% !TeX root = ../../Book1.tex
% 纲要标题：### 1.2 度量、配对与二次型
% 纲要位置：几何学计划纲要.md，第 127 行。
% 正文状态：初稿；尚未正式审核。

% 纲要内容开始
% 
% * 内积、对偶与正交投影
% * 自伴谱分解与正交标准形
% * 对称双线性型的合同分类与惯性
% * Gram 行列式、体积与退化的检测
% 
% 纲要内容结束

\section{度量、配对与二次型}
\label{g1:sec:0102}

仿射结构能说明三个点是否共线、两条直线是否平行，却不能决定一条线段有多长。把平面沿一个方向拉伸，上一节讨论的仿射关系仍然成立，圆却会变成椭圆。要分辨这样的变化，就需要在方向向量空间上再指定一种结构，用来测量长度和夹角。内积承担这一任务；更一般的对称双线性型则让我们比较正负方向、识别退化，并描述二次方程。

\ProseParagraphBreak

本节始终在实有限维向量空间上讨论。线性映射可以用矩阵表示，但同一个矩阵也可能表示双线性型；二者换基时的公式不同。先弄清对象，再讨论它的标准形，是后面几何分类中避免混淆的关键。

\subsection{内积、对偶与正交投影}
\label{g1:subsec:0102:01}

设 \(V\) 是实向量空间。我们希望长度的平方是一个二次表达式，并且非零向量的长度严格为正。为此，先规定两个向量之间的配对。

\begin{definition}[内积与欧氏结构]
	\label{g1:def:inner-product}
	实向量空间 \(V\) 上的\term[neiji]{内积}[inner product]是映射
	\[
		\langle\,\cdot\,,\,\cdot\,\rangle\colon V\times V\longrightarrow\mathbb R,
	\]
	满足对两个变量分别线性、\(\langle u,v\rangle=\langle v,u\rangle\)，并且对每个非零 \(v\in V\) 有 \(\langle v,v\rangle>0\)。最后一个条件称为\term[zhengdingxing]{正定性}[positive definiteness]。记
	\[
		\lVert v\rVert\coloneqq\sqrt{\langle v,v\rangle},
	\]
	称为内积给出的\term[fanshu]{范数}[norm]。带有内积的实有限维向量空间称为\term[oushixiangliangkongjian]{欧氏向量空间}[Euclidean vector space]。若仿射空间 \(A\) 的方向空间 \(V\) 带有内积，则称 \(A\) 带有\term[oushijiegou]{欧氏结构}[Euclidean structure]，两点的距离定义为
	\[
		d(p,q)\coloneqq\lVert q-p\rVert.
	\]
	\symidx{\(\langle u,v\rangle\)：内积}
	\symidx{\(\lVert v\rVert\)：内积范数}
\end{definition}

这里的正定性有两项直接作用。它保证长度为零的向量只有零向量，也保证长度满足三角不等式。后一项可以从一个简单的非负平方看出。若 \(v\ne0\)，则
\[
	0\leq\left\lVert u-\frac{\langle u,v\rangle}{\lVert v\rVert^2}v\right\rVert^2
	=\lVert u\rVert^2-\frac{\langle u,v\rangle^2}{\lVert v\rVert^2}.
\]
所以 \(\lvert\langle u,v\rangle\rvert\leq\lVert u\rVert\cdot\lVert v\rVert\)，\(v=0\) 时也成立。代入
\[
	\lVert u+v\rVert^2=\lVert u\rVert^2+2\langle u,v\rangle+\lVert v\rVert^2
	\leq(\lVert u\rVert+\lVert v\rVert)^2
\]
便得到三角不等式。结合上一节的点差相加公式，\(d\) 确实满足距离的三角不等式。对于两个非零向量，上述配对不等式还保证
\[
	\cos\theta=\frac{\langle u,v\rangle}{\lVert u\rVert\cdot\lVert v\rVert}
\]
确定唯一的 \(\theta\in[0,\pi]\)，这就是它们的\term[jiajiao]{夹角}[angle]。

\begin{proposition}[内积的恢复]
	\label{g1:prop:polarization}
	在实内积空间 \(V\) 上，长度确定内积，具体有
	\begin{equation}
		\label{g1:eq:inner-polarization}
		\langle u,v\rangle
		=\frac{\lVert u+v\rVert^2-\lVert u\rVert^2-\lVert v\rVert^2}{2}
		=\frac{\lVert u+v\rVert^2-\lVert u-v\rVert^2}{4}.
	\end{equation}
\end{proposition}

\begin{proof}
	由双线性和对称性展开 \(\langle u+v,u+v\rangle\) 与 \(\langle u-v,u-v\rangle\)，交叉项分别是 \(2\langle u,v\rangle\) 与 \(-2\langle u,v\rangle\)。整理即得两式。
\end{proof}

公式\eqref{g1:eq:inner-polarization}称为内积的\term[jihuahengdengshi]{极化恒等式}[polarization identity]。它说明保持所有长度的线性映射也保持所有内积。这里的“线性”暂时不可省去，因为展开时要知道映射怎样处理 \(u+v\)；下一节将解释，保持所有点间距离的映射为什么会恢复出仿射结构。

\ProseParagraphBreak

内积还使向量和线性函数联系起来。记 \(V^*\) 为 \(V\) 的\term[duioukongjian]{对偶空间}[dual space]，即所有线性映射 \(\alpha\colon V\rightarrow\mathbb R\) 组成的向量空间。\(\alpha(v)\) 是不需要内积就能定义的\term[duioupeidui]{对偶配对}[duality pairing]。给定内积以后，每个向量 \(v\) 又产生一个线性函数
\[
	v^\flat\colon w\longmapsto\langle v,w\rangle.
\]
映射 \(\flat\colon V\rightarrow V^*\) 是线性的；若 \(v^\flat=0\)，取 \(w=v\) 得 \(v=0\)，故它单射。给定 \(V\) 的一组基，逐个取坐标的线性函数构成 \(V^*\) 的一组基：任意线性函数由它在原基上的值唯一确定。因此 \(\dim V^*=\dim V\)，上述单射也是满射。于是每个 \(\alpha\in V^*\) 恰有一个向量 \(\alpha^\sharp\) 满足
\[
	\alpha(w)=\langle\alpha^\sharp,w\rangle\qquad(w\in V).
\]
\symidx{\(V^*\)：对偶空间}
\symidx{\(v^\flat,\alpha^\sharp\)：由内积给出的向量与对偶向量的对应}
对偶空间本身不依赖内积，但这项对应依赖内积，不能把它当作一个无需选择的相等关系。

\begin{example}[同一个线性函数对应不同的向量]
	\label{g1:exm:metric-dual}
	在 \(\mathbb R^2\) 上取
	\[
		\langle x,y\rangle=x_1y_1+4x_2y_2,\qquad
		\alpha(x)=x_1+2x_2.
	\]
	由 \(\alpha(x)=\langle(1,1/2),x\rangle\) 得 \(\alpha^\sharp=(1,1/2)\)。若改用通常的点积，则同一个 \(\alpha\) 对应 \((1,2)\)。更一般地，若基 \(b_1,\ldots,b_n\) 下的内积矩阵为 \(G=(\langle b_i,b_j\rangle)\)，向量 \(v\) 的坐标列为 \(c\)，则 \(v^\flat\) 的系数列是 \(Gc\)，因为 \(\langle v,w\rangle=c^{\mathsf T}Gx\)。把向量坐标与线性函数的系数直接等同，只有在 \(G=I\) 时才正确。
\end{example}

在内积为零时，称两个向量\term[zhengjiao]{正交}[orthogonal]，记 \(u\perp v\)。一组基 \(e_1,\ldots,e_n\) 若满足 \(\langle e_i,e_j\rangle=\delta_{ij}\)，便称为\term[biaozhunzhengjiaoji]{标准正交基}[orthonormal basis]。标准正交基把长度公式化成平方和；它的存在应当从内积公理推出，而不应预先藏在“欧氏”这个名称中。

\begin{proposition}[保持逐步张成空间的正交化]
	\label{g1:prop:orthonormalization}
	设 \(v_1,\ldots,v_k\) 是实有限维内积空间 \(V\) 中的线性无关向量。则存在标准正交向量组 \(e_1,\ldots,e_k\)，使得对每个 \(1\leq j\leq k\)，都有
	\[
		\operatorname{span}(e_1,\ldots,e_j)
		=\operatorname{span}(v_1,\ldots,v_j).
	\]
	特别地，\(V\) 有标准正交基，任意子空间的标准正交基都能扩充为 \(V\) 的标准正交基。
\end{proposition}

\begin{proof}
	构造的想法是从 \(v_j\) 中减去已经处理的方向，只保留新增加的方向。令
	\[
		w_1=v_1,\qquad e_1=\frac{w_1}{\lVert w_1\rVert}.
	\]
	若已经构造了前 \(j-1\) 个向量，定义
	\begin{equation}
		\label{g1:eq:orthogonalization}
		w_j=v_j-\sum_{i=1}^{j-1}\langle v_j,e_i\rangle e_i,
		\qquad e_j=\frac{w_j}{\lVert w_j\rVert}.
	\end{equation}
	对每个 \(\ell<j\)，标准正交性给出
	\[
		\langle w_j,e_\ell\rangle
		=\langle v_j,e_\ell\rangle
		-\sum_{i<j}\langle v_j,e_i\rangle\delta_{i\ell}=0.
	\]
	若 \(w_j=0\)，则 \(v_j\) 属于前 \(j-1\) 个 \(e_i\) 的张成空间，也就是前 \(j-1\) 个 \(v_i\) 的张成空间，与线性无关矛盾。因此除以 \(\lVert w_j\rVert\) 合法。式\eqref{g1:eq:orthogonalization}又说明 \(e_j\) 可由前 \(j\) 个 \(v_i\) 表出，而 \(v_j\) 可由前 \(j\) 个 \(e_i\) 表出，所以逐步张成空间相同。归纳完成构造。

	取 \(k=\dim V\) 便得到标准正交基。若从某个子空间的标准正交基出发，先把它扩充为 \(V\) 的一组普通基，再作上述构造，原有标准正交向量不变，因此得到所需扩充。
\end{proof}

这一递推构造称为\term[GramSchmidtzhengjiaohua]{Gram--Schmidt 正交化}[Gram--Schmidt orthogonalization]。\footnote{Gram（约尔根·佩德森·格拉姆，Jørgen Pedersen Gram，1850--1916），丹麦数学家，研究正交化、概率与数值方法；Schmidt（埃尔哈德·施密特，Erhard Schmidt，1876--1959），德国数学家，研究积分方程与无穷维内积空间。}输入是一组线性无关向量，输出是一组保留逐步张成空间的标准正交向量。它也说明正定性为何不可缺少：若允许非零向量的平方长度为零，构造就可能在归一化时停止。

\ProseParagraphBreak

现在考虑一个几何问题：给定向量 \(x\)，如何在子空间 \(W\) 中找到离它最近的向量？正交化提示我们，应当减去 \(x\) 在 \(W\) 的各个标准正交方向上的分量。记
\[
	W^\perp\coloneqq\{\,z\in V:\langle z,w\rangle=0\text{ 对所有 }w\in W\,\},
\]
称为 \(W\) 的\term[zhengjiaobu]{正交补}[orthogonal complement]。
\symidx{\(W^\perp\)：正交补}

\begin{proposition}[正交投影与最小距离]
	\label{g1:prop:projection}
	设 \(V\) 是实有限维内积空间，\(W\subseteq V\) 是线性子空间。每个 \(x\in V\) 都唯一地写成
	\[
		x=p+r,\qquad p\in W,\quad r\in W^\perp.
	\]
	映射 \(P_Wx=p\) 是线性的，满足 \(P_W^2=P_W\)、\(\operatorname{im}P_W=W\)、\(\ker P_W=W^\perp\)。而且对每个 \(w\in W\)，
	\begin{equation}
		\label{g1:eq:projection-distance}
		\lVert x-w\rVert^2
		=\lVert x-P_Wx\rVert^2+\lVert P_Wx-w\rVert^2.
	\end{equation}
	因此 \(P_Wx\) 是 \(W\) 中离 \(x\) 最近的唯一向量。
\end{proposition}

\begin{proof}
	取 \(W\) 的标准正交基 \(e_1,\ldots,e_k\)，令
	\[
		p=\sum_{i=1}^k\langle x,e_i\rangle e_i,\qquad r=x-p.
	\]
	逐项计算得 \(\langle r,e_j\rangle=0\)，从而 \(r\in W^\perp\)。若另有 \(x=p'+r'\)，其中 \(p'\in W\)、\(r'\in W^\perp\)，则 \(p-p'=r'-r\) 同时属于 \(W\) 和 \(W^\perp\)。这个向量与自己正交，正定性迫使它等于零。这证明唯一性，也说明构造结果不依赖所取标准正交基。

	上面的有限求和公式给出线性性。对 \(w\in W\) 有 \(P_Ww=w\)，对 \(r\in W^\perp\) 有 \(P_Wr=0\)，故幂等性、像与核的断言都成立。最后，\(x-w=(x-p)+(p-w)\) 是正交的两项之和，平方范数中的交叉项为零，得到式\eqref{g1:eq:projection-distance}。第二项非负，且恰在 \(w=p\) 时为零，故最小值在 \(p\) 处唯一取得。
\end{proof}

\(P_W\) 称为到 \(W\) 的\term[zhengjiaotouying]{正交投影}[orthogonal projection]。由分解可见
\[
	V=W\oplus W^\perp,\qquad \dim W^\perp=\dim V-\dim W,
	\qquad (W^\perp)^\perp=W.
\]
最后一式可由包含 \(W\subseteq(W^\perp)^\perp\) 和相同维数得到。\cref{g1:fig:orthogonal-projection}把距离恒等式画成一个直角三角形：换一个候选点 \(w\)，新增的平方距离正是沿子空间移动所付出的 \(\lVert P_Wx-w\rVert^2\)。
\symidx{\(P_W\)：到子空间 \(W\) 的正交投影}

\begin{figure}[htbp]
	\centering
	% !TeX root = ../Book1.tex
% 可编辑图源；由 Book1/Figures/OrthogonalProjection.tex 读入。
\begin{tikzpicture}[scale=1.0]
	\draw[->,thick] (-0.6,0)--(6.1,0) node[right]{\(W\)};
	\coordinate (O) at (0,0);
	\coordinate (P) at (3.6,0);
	\coordinate (X) at (3.6,2.5);
	\coordinate (Z) at (1.3,0);
	\draw[->,thick,color=GeometryBlue] (O)--(X);
	\draw[dashed,thick] (P)--(X);
	\draw[thick] (Z)--(X);
	\draw (3.35,0)--(3.35,.25)--(3.6,.25);
	\fill (O) circle (1.4pt) node[below]{\(0\)};
	\fill (P) circle (1.4pt) node[below]{\(P_Wx\)};
	\fill (X) circle (1.4pt) node[above]{\(x\)};
	\fill (Z) circle (1.4pt) node[below]{\(w\)};
	\node[right] at (3.6,1.25){\(x-P_Wx\)};
\end{tikzpicture}

	\caption[正交投影与最小距离]{正交投影与最小距离。图中以一维子空间示意，\(x-P_Wx\) 垂直于整个 \(W\)，因而也垂直于从 \(P_Wx\) 到任意 \(w\in W\) 的位移。}
	\label{g1:fig:orthogonal-projection}
\end{figure}

若 \(F=a+W\) 是欧氏仿射空间中的非空仿射子空间，最近点则是
\[
	\pi_F(x)=a+P_W(x-a).
\]
这是把原点临时移到 \(a\) 后应用\cref{g1:prop:projection}。若改用 \(a'=a+w_0\)，其中 \(w_0\in W\)，由 \(P_Ww_0=w_0\) 可得同一结果，故该公式不依赖基点选择。它刻画的最近点满足 \(x-\pi_F(x)\perp W\)，而非一定与位置向量 \(\pi_F(x)\) 正交。

\begin{example}[到一个仿射平面的距离]
	\label{g1:exm:plane-projection}
	在带通常点积的 \(\mathbb R^3\) 中，令 \(F=\{\,(x_1,x_2,x_3):x_1+x_2+x_3=1\,\}\)。其方向空间是 \(W=\{\,v:v_1+v_2+v_3=0\,\}\)，法向量 \(n=(1,1,1)\) 与 \(W\) 正交。因此原点的最近点必须形如 \(cn\)，再由 \(cn\in F\) 得 \(3c=1\)。最近点是 \((1/3,1/3,1/3)\)，距离为 \(1/\sqrt3\)。更一般地，若 \(n\ne0\)，到超平面 \(\langle n,y\rangle=b\) 的最近点为
	\[
		x-\frac{\langle n,x\rangle-b}{\lVert n\rVert^2}n,
	\]
	因为这个点在超平面上，而它与 \(x\) 的差沿着法向方向。
\end{example}

\subsection{自伴谱分解与正交标准形}
\label{g1:subsec:0102:02}

在标准正交坐标下，内积本身是平方和；若另有一个二次表达式，通常还会含有交叉项。例如 \(2x^2+2xy+2y^2\) 的等值线沿原来的坐标轴并不对称。是否可以只旋转或反射坐标轴，就把这些交叉项消去？答案来自自伴算子的特征方向。

\ProseParagraphBreak

给定实有限维内积空间上的线性算子 \(T\)，固定 \(v\) 后，\(u\mapsto\langle Tu,v\rangle\) 是线性函数。由本节已经证明的向量与对偶向量的对应，存在唯一向量 \(T^*v\)，使
\[
	\langle Tu,v\rangle=\langle u,T^*v\rangle\qquad(u,v\in V).
\]
唯一性同时说明 \(v\mapsto T^*v\) 线性。称 \(T^*\) 为 \(T\) 的\term[banshuisuanzi]{伴随算子}[adjoint operator]。在标准正交基下，它的矩阵是 \(T\) 的矩阵的转置；从定义连续使用两次配对等式还得到 \((AB)^*=B^*A^*\)、\((T^*)^*=T\)。例如 \(\langle ABu,v\rangle=\langle Bu,A^*v\rangle=\langle u,B^*A^*v\rangle\)，这说明乘积伴随为何颠倒次序。
\symidx{\(T^*\)：内积伴随算子}

\begin{definition}[自伴算子]
	\label{g1:def:self-adjoint}
	实内积空间 \(V\) 上的线性算子 \(T\colon V\rightarrow V\) 称为\term[zibansuanzi]{自伴算子}[self-adjoint operator]，若
	\[
		\langle Tu,v\rangle=\langle u,Tv\rangle\qquad(u,v\in V).
	\]
	若存在非零向量 \(v\) 与实数 \(\lambda\) 满足 \(Tv=\lambda v\)，则称 \(\lambda\) 为 \(T\) 的\term[tezhengzhi]{特征值}[eigenvalue]，\(v\) 为相应\term[tezhengxiangliang]{特征向量}[eigenvector]，并记\term[tezhengzikongjian]{特征子空间}[eigenspace]为
	\[
		E_\lambda(T)\coloneqq\ker(T-\lambda I).
	\]
	\symidx{\(E_\lambda(T)\)：特征子空间}
\end{definition}

若内积矩阵为 \(G\)，算子矩阵为 \(A\)，自伴条件写成 \(A^{\mathsf T}G=GA\)。只有在标准正交基下，即 \(G=I\) 时，它才简化为 \(A^{\mathsf T}=A\)。因此“矩阵对称”是自伴性在特定类型坐标中的表达。

\begin{theorem}[实自伴谱定理]
	\label{g1:thm:spectral}
	设 \(V\) 是 \(n\) 维实内积空间，\(T\colon V\rightarrow V\) 是自伴线性算子。则 \(V\) 有由 \(T\) 的特征向量组成的标准正交基。在这样的基下，\(T\) 的矩阵为
	\[
		\operatorname{diag}(\lambda_1,\ldots,\lambda_n),\qquad\lambda_i\in\mathbb R.
	\]
	不同特征值的特征子空间互相正交，\(V\) 是这些特征子空间的正交直和。特征值连同各自重数所组成的多重集由 \(T\) 唯一确定。
\end{theorem}

\begin{proof}
	关键是先找到一个实特征方向，再证明它的正交补仍由 \(T\) 保持，这样就能降低维数。若 \(n=0\)，结论由空基成立。设 \(n\geq1\)。由\cref{g1:prop:orthonormalization}，可以先取任意标准正交坐标。在这些坐标下，单位球面 \(\{\,x:\lVert x\rVert=1\,\}\) 是 \(\mathbb R^n\) 中的闭有界集合，因而紧；连续实函数在非空紧集上取得最大值。这里使用的是有限维实分析中的紧性与极值定理，有限维和正定性共同保证这个球面是通常的紧球面。

	函数 \(f(x)=\langle Tx,x\rangle\) 连续，所以在某个单位向量 \(u\) 处取得球面最大值，记为 \(\lambda=\langle Tu,u\rangle\)。对任意 \(v\perp u\)，曲线
	\[
		\gamma(t)=\frac{u+tv}{\sqrt{1+t^2\lVert v\rVert^2}}
	\]
	始终在单位球面上，并经过 \(u\)。因此 \(f(\gamma(t))\) 在 \(t=0\) 有极大值，其导数为零。利用自伴性直接展开：
	\[
		f(\gamma(t))
		=\frac{\lambda+2t\langle Tu,v\rangle+t^2\langle Tv,v\rangle}
		{1+t^2\lVert v\rVert^2}.
	\]
	在 \(t=0\) 求导得 \(2\langle Tu,v\rangle=0\)。所以 \(Tu\) 与所有 \(v\in u^\perp\) 正交。由正交分解，\(Tu\in\mathbb Ru\)；再与 \(u\) 配对得到 \(Tu=\lambda u\)。

	若 \(v\in u^\perp\)，则
	\[
		\langle Tv,u\rangle=\langle v,Tu\rangle
		=\lambda\langle v,u\rangle=0.
	\]
	故 \(u^\perp\) 在 \(T\) 下不变，限制 \(T|_{u^\perp}\) 仍满足自伴条件。对维数 \(n-1\) 应用归纳假设，得到 \(u^\perp\) 的标准正交特征基，与 \(u\) 合起来就是所求的基。

	若 \(Tu=\lambda u\)、\(Tv=\mu v\) 且 \(\lambda\ne\mu\)，则
	\[
		\lambda\langle u,v\rangle
		=\langle Tu,v\rangle
		=\langle u,Tv\rangle
		=\mu\langle u,v\rangle,
	\]
	所以 \(u\perp v\)。已有特征基又说明这些特征子空间张成整个 \(V\)。

	最后说明唯一性的含义。若换基矩阵为可逆矩阵 \(C\)，同一算子的两个矩阵为 \(A\) 与 \(C^{-1}AC\)，于是
	\[
		\det(tI-C^{-1}AC)=\det\bigl(C^{-1}(tI-A)C\bigr)=\det(tI-A).
	\]
	在所得特征基下，这个多项式等于 \(\prod_{i=1}^n(t-\lambda_i)\)。实一次因子的根和重数唯一，因而特征值的多重集不依赖所选特征基。
\end{proof}

这个结果称为实有限维的\term[pudingli]{谱定理}[spectral theorem]。“谱”在这里就是全部特征值。定理不声称特征基唯一：每个重特征值的特征子空间内，可以重新选择任意标准正交基；即使特征值互异，单位特征向量仍可改变符号。真正不依赖这些选择的是特征子空间及其正交投影，算子可以写成
\[
	T=\sum_{\lambda}\lambda P_{E_\lambda(T)}.
\]
若 \(x=\sum_i x_i e_i\)，则
\[
	\langle Tx,x\rangle=\sum_i\lambda_i x_i^2.
\]
因此 \(\langle Tx,x\rangle>0\) 对所有 \(x\ne0\) 成立，当且仅当全部 \(\lambda_i>0\)；非负性则对应全部 \(\lambda_i\geq0\)。

\begin{corollary}[实对称矩阵的正交对角化]
	\label{g1:cor:orthogonal-diagonalization}
	对每个实对称 \(n\times n\) 矩阵 \(A\)，存在满足 \(Q^{\mathsf T}Q=I\) 的实矩阵 \(Q\)，使
	\[
		Q^{\mathsf T}AQ=\operatorname{diag}(\lambda_1,\ldots,\lambda_n).
	\]
	两个同阶实对称矩阵能由这样的正交换基相互转化，当且仅当它们的特征值多重集相同。
\end{corollary}

\begin{proof}
	把 \(A\) 看作通常点积空间 \(\mathbb R^n\) 上的算子，它因 \(A^{\mathsf T}=A\) 而自伴。把\cref{g1:thm:spectral}所得标准正交特征基作为 \(Q\) 的各列，便有 \(Q^{\mathsf T}Q=I\) 及所示公式。正交换基也是相似变换，因此保持特征值多重集。反过来，若两个矩阵具有相同特征值多重集，调整特征基顺序，使两者的对角矩阵相同，复合这两个正交换基便得到所求变换。
\end{proof}

\begin{example}[椭圆的轴从哪里来]
	\label{g1:exm:ellipse-spectral}
	对称矩阵
	\[
		A=\begin{pmatrix}2&1\\1&2\end{pmatrix}
	\]
	分别在单位向量 \(e_+=(1,1)/\sqrt2\)、\(e_-=(1,-1)/\sqrt2\) 上取特征值 \(3\)、\(1\)。若 \((x,y)=u e_++v e_-\)，则
	\[
		2x^2+2xy+2y^2=3u^2+v^2.
	\]
	所以等值线 \(2x^2+2xy+2y^2=1\) 的两条主轴沿 \(e_+\)、\(e_-\)，对应半轴长分别为 \(1/\sqrt3\)、\(1\)。较大的特征值对应较短的半轴，因为沿该方向走同样的距离，二次表达式增长得更快。这里换成 \((u,v)\) 没有改变长度，只改变了标准正交坐标的方向。
\end{example}

\subsection{对称双线性型的合同分类与惯性}
\label{g1:subsec:0102:03}

正定性适合测量距离，但许多二次方程具有正负号，例如 \(x^2-y^2=1\)。这时仍有对称双线性配对，却不能把其对角值的平方根当作长度。我们因此把对称性与正定性分开。

\begin{definition}[对称双线性型与二次型]
	\label{g1:def:symmetric-form}
	设 \(V\) 是实向量空间。映射 \(B\colon V\times V\rightarrow\mathbb R\) 若对每个变量分别线性，并满足 \(B(u,v)=B(v,u)\)，称为\term[duichengshuangxianxingxing]{对称双线性型}[symmetric bilinear form]。其对应的\term[ercixing]{二次型}[quadratic form]是
	\[
		q(v)\coloneqq B(v,v).
	\]
	其\term[genji]{根基}[radical]定义为
	\[
		\operatorname{rad}B\coloneqq\{\,v\in V:B(v,w)=0\text{ 对所有 }w\in V\,\}.
	\]
	若 \(\operatorname{rad}B=\{0\}\)，称 \(B\) \term[feituihua]{非退化}[nondegenerate]；否则称其退化。
	\symidx{\(\operatorname{rad}B\)：双线性型的根基}
\end{definition}

实数域上，\(q\) 确定 \(B\)，因为
\[
	B(u,v)=\frac{q(u+v)-q(u)-q(v)}2.
\]
这里与式\eqref{g1:eq:inner-polarization}有相同的展开过程，完全不需要正定性。若取基 \(b_1,\ldots,b_n\)，令 \(A_{ij}=B(b_i,b_j)\)，则 \(A\) 对称，且坐标列为 \(x,y\) 的两个向量满足
\[
	B(u,v)=x^{\mathsf T}Ay,\qquad q(u)=x^{\mathsf T}Ax.
\]
这也说明二次多项式中的交叉项系数是相应非对角矩阵元的两倍。

\ProseParagraphBreak

若新基的各向量在旧基下组成列矩阵 \(C\)，则同一向量的旧坐标与新坐标满足 \(x=Cx'\)。因此同一个双线性型在新基下的矩阵是
\[
	A'=C^{\mathsf T}AC.
\]
这称为\term[hetong]{合同}[congruence]变换。相比之下，同一个线性算子的矩阵变为 \(C^{-1}AC\)，称为\term[xiangsi]{相似}[similarity]变换。\cref{g1:tab:similarity-congruence}列出区别。只在 \(C\) 正交、即 \(C^{-1}=C^{\mathsf T}\) 时，两种公式才重合。

\begin{table}[htbp]
	\centering
	\caption{相似变换与合同变换}
	\label{g1:tab:similarity-congruence}
	\begin{tabularx}{\linewidth}{@{}lXX@{}}
		\toprule
		& 线性算子 & 对称双线性型 \\
		\midrule
		作用对象 & 一个输入向量与一个输出向量 & 两个输入向量与一个实数值 \\
		换基公式 & \(C^{-1}AC\) & \(C^{\mathsf T}AC\) \\
		一般换基所保留 & 特征多项式、特征值 & 正负方向的最大维数、根基维数 \\
		\bottomrule
	\end{tabularx}
\end{table}

为了运用刚证得的谱定理，可以在 \(V\) 上临时选择一个辅助内积。由向量与对偶空间的对应，对每个 \(u\) 恰有一个 \(Tu\) 使
\[
	B(u,v)=\langle Tu,v\rangle\qquad(v\in V).
\]
配对的线性性和表示的唯一性说明 \(T\) 是线性算子，\(B\) 的对称性说明 \(T\) 自伴。谱定理于是把 \(B\) 化成实对角形。辅助内积只是获得某组基的工具；接下来的唯一性证明会把最终分类量刻画成 \(B\) 自己的性质。

\begin{theorem}[实对称双线性型的惯性定理]
	\label{g1:thm:inertia}
	设 \(V\) 是 \(n\) 维实向量空间，\(B\) 是 \(V\) 上的对称双线性型。则存在一组基，使 \(B\) 的矩阵为
	\[
		\operatorname{diag}(I_p,-I_q,0_r),\qquad p+q+r=n.
	\]
	三元组 \((p,q,r)\) 唯一。其中 \(p\) 是使 \(B|_{U\times U}\) 正定的子空间 \(U\) 的最大维数，\(q\) 是使 \(-B|_{U\times U}\) 正定的子空间 \(U\) 的最大维数，\(r=\dim\operatorname{rad}B\)。两个同维实对称双线性型合同，当且仅当这三个整数相同。
\end{theorem}

\begin{proof}
	先证明标准形存在。任取辅助内积，将 \(B\) 写成 \(\langle Tu,v\rangle\)。由\cref{g1:thm:spectral}，存在标准正交特征基 \(e_1,\ldots,e_n\)，其中
	\[
		B(e_i,e_j)=\lambda_i\delta_{ij}.
	\]
	对每个非零 \(\lambda_i\)，把 \(e_i\) 改为 \(e_i/\sqrt{\lvert\lambda_i\rvert}\)，对应的对角元就变成 \(1\) 或 \(-1\)；零特征值对应的向量保持不变。把正、负、零三组重新排序便得到所示标准形。这次缩放通常不是正交换基，所以它可以改变特征值的绝对大小。

	再证明计数与基无关。设所得三组向量张成的空间分别为 \(V_+\)、\(V_-\)、\(V_0\)。显然 \(B\) 在 \(V_+\) 上正定，故正定子空间的最大维数至少为 \(p\)。若 \(U\) 是任意正定子空间，则
	\[
		U\cap(V_-\oplus V_0)=\{0\},
	\]
	因为右边空间中的每个向量都满足 \(B(v,v)\leq0\)，而 \(U\) 中每个非零向量满足 \(B(v,v)>0\)。将 \(U\) 沿 \(V_-\oplus V_0\) 投影到 \(V_+\)，上述交集条件说明该线性映射单射，故 \(\dim U\leq\dim V_+=p\)。因此 \(p\) 恰等于所有正定子空间的最大维数。对 \(-B\) 应用同一论证得负定方向的最大维数为 \(q\)。

	在标准形坐标中，\(B(v,w)=0\) 对所有 \(w\) 成立，当且仅当 \(v\) 的前 \(p+q\) 个坐标都为零。因此 \(\operatorname{rad}B=V_0\)，其维数是 \(r\)。这些描述只涉及 \(B\)，没有涉及所选基或辅助内积，故三元组唯一。合同变换是同一配对的换基，必保持这三个数；反过来，三元组相同的两个型都能化为同一个标准矩阵，复合相应换基就得到二者的合同关系。
\end{proof}

三元组 \((p,q,r)\) 称为 \(B\) 的\term[guanxing]{惯性}[inertia]；\(p\)、\(q\) 分别称\term[zhengguanxingzhishu]{正惯性指数}[positive index of inertia]与\term[fuguanxingzhishu]{负惯性指数}[negative index of inertia]。\(p+q\) 是 \(B\) 的秩，\(r\) 是零化度。

\begin{example}[零值方向不等于退化方向]
	\label{g1:exm:isotropic-radical}
	在 \(\mathbb R^3\) 上考虑
	\[
		q(x,y,z)=x^2+4xy+y^2=(x+2y)^2-3y^2.
	\]
	取新坐标 \(u=x+2y\)、\(v=\sqrt3\,y\)、\(w=z\)，得到 \(q=u^2-v^2\)，惯性为 \((1,1,1)\)，根基为 \(z\) 轴。另一方面，在 \((u,v,w)\) 坐标中，向量 \((1,1,0)\) 满足 \(q=0\)，却不在根基中，因为它与 \((1,0,0)\) 的配对为 \(1\)。非零而满足 \(q(v)=0\) 的向量称为\term[mixiangxiangliang]{迷向向量}[isotropic vector]。迷向性只检查向量与自己的配对，根基则要求它与所有向量的配对都为零。这一差别也说明，不定二次型不能直接给出距离。
\end{example}

\begin{example}[相同惯性仍可有不同轴长]
	\label{g1:exm:congruence-vs-orthogonal}
	矩阵 \(I_2\) 与 \(\operatorname{diag}(1,4)\) 都有惯性 \((2,0,0)\)，并且
	\[
		\begin{pmatrix}1&0\\0&1/2\end{pmatrix}^{\!\mathsf T}
		\begin{pmatrix}1&0\\0&4\end{pmatrix}
		\begin{pmatrix}1&0\\0&1/2\end{pmatrix}=I_2.
	\]
	它们因而合同，却不能由正交换基相互转化，因为特征值多重集分别为 \(\{1,1\}\) 和 \(\{1,4\}\)。几何上，一个坐标方向的缩放可以把椭圆变成圆，保持长度的坐标变换却不能做到。允许的变换越多，分类留下的信息就可能越少。
\end{example}

\subsection{Gram 行列式、体积与退化的检测}
\label{g1:subsec:0102:04}

测量一条边只需要它的长度；测量一个由多条边张成的平行四边形或平行多面体，还需要知道这些边之间是否接近平行。把所有两两内积放进一个矩阵，正好同时保存长度与夹角的信息。

\begin{definition}[Gram 矩阵]
	\label{g1:def:gram-matrix}
	设 \(v_1,\ldots,v_k\) 是实内积空间中的向量，它们的\term[Gramjuzhen]{Gram 矩阵}[Gram matrix]为
	\[
		\Gamma(v_1,\ldots,v_k)\coloneqq
		\bigl(\langle v_i,v_j\rangle\bigr)_{1\leq i,j\leq k}.
	\]
	其行列式称为\term[Gramhanglieshi]{Gram 行列式}[Gram determinant]。
	\symidx{\(\Gamma(v_1,\ldots,v_k)\)：Gram 矩阵}
\end{definition}

对任意列向量 \(c=(c_1,\ldots,c_k)^{\mathsf T}\)，双线性给出
\begin{equation}
	\label{g1:eq:gram-quadratic}
	c^{\mathsf T}\Gamma c
	=\left\lVert\sum_{i=1}^k c_i v_i\right\rVert^2.
\end{equation}
因此 \(\Gamma\) 总是\term[banzhengding]{半正定的}[positive semidefinite]，即对应二次型处处非负；只有当这些向量线性无关时，它才正定。这个等式已经预示：Gram 行列式为零应当对应边向量的线性相关，而它的正平方根应当测量所张成的体积。

\ProseParagraphBreak

先明确体积的含义。对线性无关的 \(v_1,\ldots,v_k\)，记 \(W=\operatorname{span}(v_1,\ldots,v_k)\)。取 \(W\) 的标准正交基，把这些向量写成一个 \(k\times k\) 坐标矩阵 \(C\) 的各列。它们张成的平行多面体为
\[
	\mathcal P(v_1,\ldots,v_k)
	=\left\{\,\sum_{i=1}^k t_iv_i:0\leq t_i\leq1\,\right\},
\]
其 \(k\) 维欧氏体积定义为 \(\lvert\det C\rvert\)。这个定义采用了单位正交立方体体积为 \(1\) 的归一化：有向体积对各边线性、交换两边变号，且在所选的标准正交基上的值为 \(1\)，按各列的坐标展开就得到行列式公式。取绝对值后不再记录定向。

\ProseParagraphBreak

若换一个标准正交基，坐标矩阵变为 \(Q^{\mathsf T}C\)，其中 \(Q^{\mathsf T}Q=I\)。取行列式可得 \((\det Q)^2=1\)，所以绝对行列式不变。体积因此不依赖所选标准正交坐标。若 \(v_1,\ldots,v_k\) 线性相关，我们把它们的 \(k\) 维体积定义为零：此时多面体实际上落在低于 \(k\) 维的空间内。这个约定不否认它可能仍有非零的较低维体积。\(k=0\) 时零维平行多面体 \(\{0\}\) 的体积约定为 \(1\)。

\begin{proposition}[Gram 行列式与体积]
	\label{g1:prop:gram-volume}
	设 \(V\) 是 \(n\) 维实内积空间，\(v_1,\ldots,v_k\in V\)。则
	\[
		\operatorname{Vol}_k\bigl(\mathcal P(v_1,\ldots,v_k)\bigr)^2
		=\det\Gamma(v_1,\ldots,v_k).
	\]
	Gram 矩阵的秩等于 \(\operatorname{span}(v_1,\ldots,v_k)\) 的维数。特别地，当 \(k\geq1\) 时，Gram 行列式为正当且仅当各向量线性无关，为零当且仅当各向量线性相关；若 \(k>n\)，它必为零。
\end{proposition}

\begin{proof}
	令 \(L\colon\mathbb R^k\rightarrow V\) 为 \(L(c)=\sum_i c_iv_i\)。若 \(L(c)=0\)，则 \(\Gamma c\) 的第 \(j\) 个分量 \(\langle v_j,L(c)\rangle\) 为零，故 \(\ker L\subseteq\ker\Gamma\)。反之，若 \(\Gamma c=0\)，式\eqref{g1:eq:gram-quadratic}给出 \(\lVert L(c)\rVert^2=0\)，由正定性知 \(L(c)=0\)。因此两者的核相同。对同一个定义域应用秩与零化度公式，得到
	\[
		\operatorname{rank}\Gamma=\operatorname{rank}L
		=\dim\operatorname{span}(v_1,\ldots,v_k).
	\]
	若各向量线性相关，\(\Gamma\) 不可逆，行列式和所定义的 \(k\) 维体积都为零。

	若各向量线性无关，在它们张成的 \(k\) 维空间中取标准正交基，记坐标矩阵为 \(C\)。逐项计算有 \(\Gamma=C^{\mathsf T}C\)，所以
	\[
		\det\Gamma=\det(C^{\mathsf T})\det C=(\det C)^2
		=\operatorname{Vol}_k(\mathcal P)^2>0.
	\]
	线性无关时行列式为正，与线性相关时的零值一起给出了退化判据。\(k=0\) 时空矩阵的行列式约定为 \(1\)，与体积约定相容。
\end{proof}

体积公式还可以从“底乘高”理解。对线性无关的边向量作式\eqref{g1:eq:orthogonalization}的正交化。记
\[
	h_j=\lVert w_j\rVert
	=d\bigl(v_j,\operatorname{span}(v_1,\ldots,v_{j-1})\bigr),
\]
其中第一步的张成空间是 \(\{0\}\)。由\cref{g1:prop:projection}，\(w_j\) 正是减去正交投影后的剩余向量。在所得标准正交基中，第 \(j\) 列只在前 \(j\) 行有非零坐标，对角元为 \(h_j\)。因此坐标矩阵上三角，得到
\[
	\operatorname{Vol}_k(\mathcal P)=h_1h_2\cdots h_k.
\]
每增加一条边，体积乘上它离已有边所张成空间的高度。任一高度为零都意味着并未增加一个新维度。

\begin{example}[三维空间中的面积与体积]
	\label{g1:exm:gram-volume}
	在通常的 \(\mathbb R^3\) 中取
	\[
		v_1=(1,1,0),\qquad v_2=(1,0,1),\qquad v_3=(0,1,1).
	\]
	前两条边的 Gram 矩阵为
	\[
		\begin{pmatrix}2&1\\1&2\end{pmatrix},
	\]
	故平行四边形面积为 \(\sqrt3\)。同一数也可由底乘高得到：底长是 \(\sqrt2\)，而
	\[
		\left\lVert v_2-\frac{\langle v_2,v_1\rangle}{\lVert v_1\rVert^2}v_1\right\rVert^2
		=2-\frac12=\frac32.
	\]
	三条边的 Gram 矩阵为
	\[
		\begin{pmatrix}2&1&1\\1&2&1\\1&1&2\end{pmatrix},
	\]
	其行列式为 \(4\)，所以平行六面体体积为 \(2\)。各边长都是 \(\sqrt2\)，但体积不是 \((\sqrt2)^3\)，因为三条边并不两两正交。
\end{example}

\begin{example}[长度不消失而面积退化]
	\label{g1:exm:gram-degeneration}
	在 \(\mathbb R^2\) 中令 \(v_1=(1,0)\)、\(v_2(t)=(1,t)\)。它们的长度在 \(t\rightarrow0\) 时都趋于 \(1\)，但
	\[
		\Gamma(t)=\begin{pmatrix}1&1\\1&1+t^2\end{pmatrix},
		\qquad \det\Gamma(t)=t^2,
	\]
	所以面积为 \(\lvert t\rvert\)，趋于零。退化发生在两个方向逐渐重合，而不是某条边变短。对任意两个非零向量，Gram 公式也写成
	\[
		\operatorname{Area}^2
		=\lVert v_1\rVert^2\lVert v_2\rVert^2-\langle v_1,v_2\rangle^2
		=\lVert v_1\rVert^2\lVert v_2\rVert^2\sin^2\theta,
	\]
	这正是熟悉的“边长乘边长再乘夹角正弦”。
\end{example}

最后要保留正定性这个条件。若把 Gram 矩阵中的内积换成任意对称双线性型，行列式为零便未必表示线性相关。例如 \(B((x,y),(x',y'))=xx'-yy'\) 非退化，而非零向量 \(v=(1,1)\) 所组成的一元向量组线性无关，其配对矩阵却是 \((0)\)。\cref{g1:prop:gram-volume}证明中从平方长度为零推出向量为零的一步，此时已经失效。

% !TeX root = ../../Book1.tex
% 纲要标题：### 1.3 Euclid 等距变换的结构定理
% 纲要位置：几何学计划纲要.md，第 134 行。
% 正文状态：初稿；尚未正式审核。

% 纲要内容开始
% 
% * 由距离保持恢复仿射性
% 
%   先用中点及极化恢复线性结构，再分离正交部分和平移部分。
% 
% * 正交算子的实标准分块
% * 不动点、螺旋运动与平移分量
% * 反射生成与低维分类
% 
% 纲要内容结束

\section{Euclid 等距变换的结构定理}
\label{g1:sec:0103}

一张刚性平面图形可以平移、旋转，也可以翻转后与原图形重合。这些动作都保持两点间的距离，但只从“距离保持”出发，是否还可能有形状更复杂的变换？关键是用距离唯一确定中点：保持距离便保持中点，连续性再把中点关系推广为任意仿射组合。得到仿射形式后，正交部分控制旋转与反射，平移部分则由不动方向决定哪些分量可以通过换原点消去。

\subsection{由距离保持恢复仿射性}
\label{g1:subsec:0103:01}

在一个有限维实仿射空间的平移向量空间上给定正定内积，就得到了\term[oushikongjian]{Euclid 空间}[Euclidean space]。其距离定义为 \(d(p,q)=\lVert q-p\rVert\)。选取原点和标准正交方向基后，可以将它写成 \(\mathbb R^n\)，内积成为通常的点积；更换原点不会改变点差。以下先在这种坐标中论证，再把结果解释为不依赖原点的几何结构。

\begin{definition}
	\label{g1:def:isometry}
	设 \(E\) 是有限维 Euclid 空间。双射 \(f\colon E\rightarrow E\) 若对所有 \(p,q\in E\) 满足
	\[
		d(f(p),f(q))=d(p,q),
	\]
	就称为 \(E\) 的\term[dengjubianhuan]{等距变换}[isometry]。若暂不要求双射，只要求上述等式，称 \(f\) 为距离保持映射。
\end{definition}

距离保持首先给出连续性：两像点之间的距离恰好等于原距离。更关键的是，它还给出了中点。对任意 \(a,b,z\in\mathbb R^n\)，展开内积可得
\[
	\lVert z-a\rVert^2+\lVert z-b\rVert^2
	=2\left\lVert z-\frac{a+b}{2}\right\rVert^2
	+\frac12\lVert a-b\rVert^2.
\]
因此，同时满足
\[
	\lVert z-a\rVert=\lVert z-b\rVert=\frac12\lVert a-b\rVert
\]
的点只有 \((a+b)/2\)。这里既要求两段等长，也要求其长度恰好是端点距离的一半；当 \(a\ne b\) 时，仅有“到两端等距”会得到整个中垂超平面，不能确定中点。

\ProseParagraphBreak

若实有限维内积空间上的线性算子 \(Q\) 对所有向量都满足 \(\langle Qu,Qv\rangle=\langle u,v\rangle\)，就称其为\term[zhengjiaosuanzi]{正交算子}[orthogonal operator]。由伴随的定义，这等价于 \(Q^*Q=I\)。这样的算子单射，因定义域和值域同维而可逆，故 \(Q^*=Q^{-1}\)。在标准正交基下，其矩阵满足 \(Q^{\mathsf T}Q=I\)。复合与求逆都保持这一性质，所有这样的 \(n\) 阶矩阵组成\term[zhengjiaoqun]{正交群}[orthogonal group]，记为 \(O(n)\)。
\symidx{\(O(n)\)：实正交群}

\begin{theorem}[等距变换的仿射形式]
	\label{g1:thm:isometry-affine}
	设 \(n\geq 1\)，且 \(f\colon\mathbb R^n\rightarrow\mathbb R^n\) 对所有 \(x,y\) 都满足 \(\lVert f(x)-f(y)\rVert=\lVert x-y\rVert\)。那么存在唯一的向量 \(t\in\mathbb R^n\) 和唯一的正交线性算子 \(Q\)，使
	\[
		f(x)=Qx+t.
	\]
	特别地，\(f\) 自动为双射。反之，每个这种形式的映射都是等距变换。
\end{theorem}

\begin{proof}
	先证明中点保持，再利用连续性把整数、有理数的运算推广到实数。若 \(m=(x+y)/2\)，则距离保持说明 \(f(m)\) 到 \(f(x),f(y)\) 的距离均为 \(\lVert f(x)-f(y)\rVert/2\)。由刚才证明的中点唯一性，
	\[
		f\left(\frac{x+y}{2}\right)=\frac{f(x)+f(y)}2.
	\]
	这一步只涉及已知像点，不需要预先知道 \(f\) 满射。

	令 \(t=f(0)\)，\(g(x)=f(x)-t\)，则 \(g(0)=0\)，且 \(g\) 同样保持中点与距离。取端点为 \(0,2x\)，得到 \(g(2x)=2g(x)\)，从而
	\[
		g(x+y)=2g\left(\frac{x+y}{2}\right)=g(x)+g(y).
	\]
	由加法性可得 \(g(-x)=-g(x)\)、\(g(mx)=mg(x)\) 对整数 \(m\) 成立；若 \(r=m/k\) 为有理数且 \(k>0\)，则 \(kg(rx)=g(krx)=mg(x)\)，所以 \(g(rx)=rg(x)\)。对于任意实数 \(a\)，取有理数列 \(r_j\rightarrow a\)。距离保持给出
	\[
		\lVert g(r_jx)-g(ax)\rVert=|r_j-a|\lVert x\rVert\longrightarrow 0,
	\]
	因而 \(g(ax)=ag(x)\)。这就证明了 \(g\) 为实线性算子，记为 \(Q\)。

	最后由极化公式恢复内积。将\cref{g1:prop:polarization}中的平方范数展开用于 \(x-y\)，得到
	\[
		\langle Qx,Qy\rangle
		=\frac12\bigl(\lVert Qx\rVert^2+\lVert Qy\rVert^2
		-\lVert Q(x-y)\rVert^2\bigr)
		=\langle x,y\rangle.
	\]
	所以 \(Q\) 正交。特别地，\(Qx=0\) 蕴含 \(x=0\)，有限维同维线性映射的秩与核维数公式便说明 \(Q\) 可逆；于是 \(f\) 的逆为 \(y\mapsto Q^{-1}(y-t)\)。唯一性来自 \(t=f(0)\) 和 \(Qx=f(x)-f(0)\)。反向验证只需计算 \(\lVert Qx+t-(Qy+t)\rVert=\lVert Q(x-y)\rVert=\lVert x-y\rVert\)。
\end{proof}

按照\cref{g1:def:affine-map}，上一定理正是说：Euclid 空间中只用距离提出的保持条件，已经强制映射具有仿射形式。有限维和定义域、值域同维在由单射推出满射时起作用。例如 \(x\mapsto(x,0)\) 从 \(\mathbb R^n\) 到 \(\mathbb R^{n+1}\) 保持距离，却不能覆盖整个值域。

\subsection{正交算子的实标准分块}
\label{g1:subsec:0103:02}

平移部分暂时放在一边。正交算子满足 \(Q^*Q=I\)，其中 \(Q^*\) 是内积伴随；在标准正交基下，这就是矩阵转置满足 \(Q^{\mathsf T}Q=I\)。它通常并不自伴：平面上一个转角既不是 \(0\) 也不是 \(\pi\)（模 \(2\pi\)）的旋转，就没有非零实特征向量。因此不能直接把\cref{g1:thm:spectral}应用于 \(Q\)。我们将对它的自伴部分应用该定理，余下的部分恰好给出旋转平面。

\begin{theorem}[正交算子的实分块]
	\label{g1:thm:orthogonal-blocks}
	设 \(V\) 是 \(n\) 维实内积空间，\(Q\colon V\rightarrow V\) 为正交算子。存在一组标准正交基，使其矩阵成为
	\[
		\operatorname{diag}\bigl(I_p,-I_q,R_{\theta_1},\ldots,R_{\theta_k}\bigr),
		\qquad p+q+2k=n,
	\]
	其中
	\begin{equation}
		\label{g1:eq:rotation-block}
		R_\theta=
		\begin{pmatrix}\cos\theta&-\sin\theta\\
			\sin\theta&\cos\theta\end{pmatrix},
		\qquad 0<\theta<\pi.
	\end{equation}
	整数 \(p,q\) 和角度的无序多重集 \(\{\theta_1,\ldots,\theta_k\}\) 由 \(Q\) 唯一确定；各个旋转平面的选择未必唯一。特别地，\(\det Q=(-1)^q\)。
\end{theorem}

\begin{proof}
	取自伴算子
	\[
		S=\frac{Q+Q^*}{2}.
	\]
	由\cref{g1:thm:spectral}，有限维实内积空间中的自伴算子有正交特征分解，故 \(V\) 是 \(S\) 的各特征空间 \(E_\lambda\) 的正交直和。因为 \(Q^*=Q^{-1}\)，所以 \(QS=SQ\)，每个 \(E_\lambda\) 均被 \(Q\) 保持。若 \(v\in E_\lambda\)，则
	\[
		\lambda\lVert v\rVert^2=\langle Sv,v\rangle
		=\langle Qv,v\rangle,
	\]
	由内积不等式得 \(|\lambda|\leq 1\)。对于端点，
	\[
		\lVert Qv-v\rVert^2=2(1-\lambda)\lVert v\rVert^2,
		\qquad
		\lVert Qv+v\rVert^2=2(1+\lambda)\lVert v\rVert^2,
	\]
	故 \(Q|_{E_1}=I\)，\(Q|_{E_{-1}}=-I\)。

	现在设 \(-1<\lambda<1\)，并记 \(s=\sqrt{1-\lambda^2}>0\)。在 \(E_\lambda\) 上有 \(Q+Q^*=2\lambda I\)；左乘 \(Q\) 得 \(Q^2-2\lambda Q+I=0\)。因此，去掉 \(Q\) 的标量部分并重新归一化，得到
	\[
		J=\frac{Q-\lambda I}{s},\qquad J^2=-I,
		\qquad J^*=-J,
		\qquad J^*J=I.
	\]
	这个构造的作用是把一般角度旋转还原成直角旋转：\(J^2=-I\) 表明连续做两次得到反向，而 \(J^*=-J\) 给出 \(\langle e,Je\rangle=0\)。选取任意单位向量 \(e\in E_\lambda\)，则 \(e,Je\) 是正交单位向量，其张成的平面 \(P\) 被 \(J\) 保持，并且
	\[
		Qe=\lambda e+sJe,\qquad Q(Je)=-se+\lambda Je.
	\]
	取唯一的 \(\theta\in(0,\pi)\) 使 \(\cos\theta=\lambda\)、\(\sin\theta=s\)，便得到\eqref{g1:eq:rotation-block}。

	在这个平面的正交补中可以继续拆分。若 \(w\in E_\lambda\cap P^\perp\) 且 \(z\in P\)，则 \(\langle Jw,z\rangle=-\langle w,Jz\rangle=0\)，所以 \(E_\lambda\cap P^\perp\) 也被 \(J\) 保持。在这个子空间中重复选取单位向量及其 \(J\) 像，每次使维数降低二，直至空间耗尽。若余下一维，任取非零向量仍会产生与它正交的非零 \(J\) 像，矛盾。因此 \(\dim E_\lambda\) 必为偶数，产生 \(\dim E_\lambda/2\) 个相同角度的旋转块。

	最后，\(p=\dim E_1\)、\(q=\dim E_{-1}\)，其余角度及重数由 \(S\) 的特征值和特征空间维数确定，故与基的选择无关。每个旋转块行列式为 \(1\)，其余一维块为 \(1\) 或 \(-1\)，于是得到行列式公式。
\end{proof}

这里把 \(\theta=0\) 归入两个 \(+1\) 块，把 \(\theta=\pi\) 归入两个 \(-1\) 块，因而标准形的记法不会把端点重复计算。一个平面的基向量顺序反过来时，\(R_\theta\) 会变为 \(R_{-\theta}\)；定理允许任意正交基，所以可以统一取 \(0<\theta<\pi\)。若平面已有固定定向，并且只准使用与之相容的基，旋转角的正负就仍须保留。

\begin{definition}
	\label{g1:def:direct-isometry}
	对于有限维 Euclid 空间的等距变换 \(f(x)=Qx+t\)，若 \(\det Q=1\)，称其\term[baoxiangdengjubianhuan]{保向}[orientation-preserving]；若 \(\det Q=-1\)，称其\term[fanxiangdengjubianhuan]{反向}[orientation-reversing]。
\end{definition}

这一称呼可以直接用有序基理解：在两组有序基之间，换基矩阵行列式为正时将它们视为同一定向，为负时视为相反定向。正交算子的行列式只能是 \(\pm1\)，分别保持或交换这两类基；平移不改变差向量，因而不影响定向。

\subsection{不动点、螺旋运动与平移分量}
\label{g1:subsec:0103:03}

写出 \(f(x)=Qx+t\) 后，接着考察换原点怎样改变平移部分。纯平移的向量在换原点时保持不变。若新原点的旧坐标是 \(c\)，则新坐标为 \(y=x-c\)，且
\begin{equation}
	\label{g1:eq:isometry-origin-change}
	f(c+y)-c=Qy+t-(I-Q)c.
\end{equation}
能消去哪些平移分量，取决于 \(I-Q\) 的像；特别是 \(Q\) 固定的方向，不能由 \((I-Q)c\) 产生。下面把这一障碍与可消去的部分分开。

\begin{theorem}[等距变换的平移分解]
	\label{g1:thm:isometry-normal}
	设 \(f(x)=Qx+t\) 是 \(\mathbb R^n\) 的等距变换，令
	\[
		F=\ker(I-Q),\qquad t=a+b,\qquad a\in F,\quad b\in F^\perp
	\]
	为 \(t\) 的正交分解。存在唯一的 \(c\in F^\perp\) 满足 \((I-Q)c=b\)，且对所有 \(u\in F\)、\(v\in F^\perp\)，
	\begin{equation}
		\label{g1:eq:isometry-normal}
		f(c+u+v)=c+u+a+Qv.
	\end{equation}
	变换 \(f\) 有不动点当且仅当 \(a=0\)；这时不动点集恰为 \(c+F\)。对所有 \(x\in\mathbb R^n\)，有
	\[
		\lVert f(x)-x\rVert\geq\lVert a\rVert,
	\]
	且等号恰在仿射子空间 \(c+F\) 上成立。该子空间被 \(f\) 保持，\(f\) 在其上就是沿 \(a\) 的平移。
\end{theorem}

\begin{proof}
	由\cref{g1:prop:projection}，有限维内积空间中的向量有相对于 \(F\) 的唯一正交分解。若 \(v\in F^\perp\)、\(u\in F\)，则
	\[
		\langle Qv,u\rangle=\langle v,Q^*u\rangle
		=\langle v,u\rangle=0,
	\]
	其中 \(Q^*u=Q^{-1}u=u\)，故 \(F^\perp\) 被 \(Q\) 和 \(I-Q\) 保持。在 \(F^\perp\) 上，\((I-Q)v=0\) 蕴含 \(v\in F\cap F^\perp=\{0\}\)，所以这个子空间上的线性算子 \(I-Q\) 单射，也就可逆。于是存在唯一 \(c\in F^\perp\) 使 \((I-Q)c=b\)。将 \(t=a+b\) 代入\eqref{g1:eq:isometry-origin-change}，再用 \(Qu=u\)，得到\eqref{g1:eq:isometry-normal}。

	若 \(x=c+u+v\)，位移分成两个正交部分：
	\[
		f(x)-x=a+(Q-I)v,
		\qquad
		\lVert f(x)-x\rVert^2=\lVert a\rVert^2+\lVert(Q-I)v\rVert^2.
	\]
	由于 \(I-Q\) 在 \(F^\perp\) 上可逆，第二项为零当且仅当 \(v=0\)。这既给出最小位移及其取得位置，也说明存在零位移点当且仅当 \(a=0\)，此时所有零位移点为 \(c+F\)。在 \(v=0\) 时，\eqref{g1:eq:isometry-normal}还直接给出 \(f(c+u)=c+u+a\)。
\end{proof}

向量 \(a\) 是无法通过平移原点消去的部分：由\eqref{g1:eq:isometry-origin-change}，原点改变只给 \(t\) 添加 \(F^\perp\) 中的向量，所以它在 \(F\) 上的投影不变。虽然 \(c\) 的坐标依赖起初的原点，\(c+F\) 却可以直接从 \(f\) 识别为“位移最小的点的集合”。这说明分解中的几何对象比最初写下的 \(t\) 更有意义。

\begin{example}[滑移反射及其最小位移]
	\label{g1:exm:glide-reflection}
	考虑平面变换 \(f(x,y)=(x+h,-y)\)，其中 \(h\ne0\)。它先关于 \(x\) 轴反射，再沿该轴平移 \((h,0)\)，称为\term[huayifanshe]{滑移反射}[glide reflection]。这两步可以交换次序，因为反射不改变轴方向的向量。对点 \((x,y)\)，位移为 \((h,-2y)\)，长度为 \(\sqrt{h^2+4y^2}\)，所以它没有不动点；轴上的位移最小，最小值为 \(|h|\)。\cref{g1:fig:glide-displacement}把沿轴的平移和垂直于轴的翻转分别画出。若 \(h=0\)，变换退化为普通轴反射，本文不再将其列入非平凡滑移反射。
\end{example}

\begin{figure}[htbp]
	\centering
	% !TeX root = ../Book1.tex
\begin{tikzpicture}[x=1cm,y=1cm,>=stealth]
	\draw[->] (-0.6,0)--(6.0,0) node[right] {反射轴};
	\coordinate (P) at (0.7,1.4);
	\coordinate (R) at (0.7,-1.4);
	\coordinate (F) at (4.8,-1.4);
	\fill (P) circle (1.8pt) node[above] {\(p\)};
	\fill (R) circle (1.8pt) node[below] {\(r(p)\)};
	\fill (F) circle (1.8pt) node[below] {\(f(p)\)};
	\draw[dashed,->] (P)--(R);
	\draw[->,thick] (R)--(F) node[midway,above] {\(a=(h,0)\)};
	\draw[->,thick,GeometryBlue] (P)--(F);
	\draw (0.7,0.16)--(0.86,0.16)--(0.86,0);
\end{tikzpicture}

	\caption[滑移反射的两部分]{滑移反射的两部分。先由 \(p\) 反射到 \(r(p)\)，再沿轴平移到 \(f(p)\)；虚线表示垂直于轴的位移分量，蓝色斜箭头表示总位移。轴上每一点只发生长度为 \(|h|\) 的平移。}
	\label{g1:fig:glide-displacement}
\end{figure}

在三维空间中，若 \(Q\) 固定一条直线，并在其垂直平面上旋转角 \(\theta\)，正规形为
\[
	(x,y,z)\longmapsto
	\bigl(x\cos\theta-y\sin\theta,
	x\sin\theta+y\cos\theta,z+h\bigr).
\]
当 \(\theta\not\equiv0\pmod{2\pi}\) 且 \(h\ne0\) 时，称为\term[luoxuanyundong]{螺旋运动}[screw motion]。每次作用绕同一轴转动相同角度，同时沿轴前进相同距离；旋转与轴向平移可交换。距轴为 \(r\) 的点在垂直平面内移动的弦长平方为 \(2r^2(1-\cos\theta)\)，所以它的位移长度为
\[
	\sqrt{h^2+4r^2\sin^2(\theta/2)}.
\]
最小位移为 \(|h|\)，恰在轴上取得。\(h=0\) 时是绕轴旋转；\(\theta\equiv0\) 时是纯平移或恒等变换，此时不能由变换本身选出唯一的轴。

\begin{example}[从坐标公式找螺旋轴]
	\label{g1:exm:screw-axis}
	设
	\[
		f(x,y,z)=(-y+2,x+1,z+3).
	\]
	其正交部分在 \(xy\) 平面旋转 \(\pi/2\)，固定方向为 \(F=\mathbb R(0,0,1)\)。平移向量分为 \(t=(0,0,3)+(2,1,0)\)，故 \(a=(0,0,3)\)。为找出能够消去横向平移的新原点，求解 \((I-Q)c=(2,1,0)\) 且 \(c\in F^\perp\)，即
	\[
		c_x+c_y=2,\qquad -c_x+c_y=1,\qquad c_z=0.
	\]
	得到 \(c=(1/2,3/2,0)\)。在 \((u,v,w)=(x-1/2,y-3/2,z)\) 坐标中，变换成为 \((u,v,w)\mapsto(-v,u,w+3)\)。因此螺旋轴是 \(\{(1/2,3/2,z):z\in\mathbb R\}\)，每次沿轴前进 \(3\)，没有不动点。连续作用四次时旋转恢复原状，故 \(f^4(x,y,z)=(x,y,z+12)\)。
\end{example}

\subsection{反射生成与低维分类}
\label{g1:subsec:0103:04}

旋转和平移看起来与镜面反射差别很大，但它们都可以由有限次反射合成。先考虑过原点的超平面。若 \(u\ne0\)，关于 \(u^\perp\) 的\term[chaopingmianfanshe]{超平面反射}[hyperplane reflection]为
\[
	r_u(x)=x-2\frac{\langle x,u\rangle}{\langle u,u\rangle}u.
\]
向量沿 \(u^\perp\) 的分量保持不变，沿 \(u\) 的分量反号，因此 \(r_u\) 正交，\(r_u^2=I\)，\(\det r_u=-1\)。一般仿射超平面 \(c+u^\perp\) 的反射则是 \(x\mapsto c+r_u(x-c)\)。

\begin{theorem}[反射生成]
	\label{g1:thm:reflection-generation}
	\(n\) 维实内积空间的每个正交算子都是至多 \(n\) 个过原点的超平面反射的复合；恒等算子允许用零个反射表示。进一步，\(n\) 维 Euclid 空间的每个等距变换都是至多 \(n+1\) 个仿射超平面反射的复合。
\end{theorem}

\begin{proof}
	先对维数证明线性结论。零维情形只有恒等算子。设 \(n>0\)，取单位向量 \(e\)。若 \(Qe=e\)，则 \(e^\perp\) 被 \(Q\) 保持；对这一 \(n-1\) 维空间使用归纳假设，将 \(Q|_{e^\perp}\) 写成至多 \(n-1\) 个反射。每个反射的法向量都在 \(e^\perp\)，用同一法向量在整个空间定义反射，便在 \(e\) 方向上保持恒等，且其复合正是 \(Q\)。

	若 \(Qe\ne e\)，我们先用一次反射把 \(Qe\) 送回 \(e\)，之后就能沿用上一步的降维。令 \(u=Qe-e\)，则
	\[
		\lVert u\rVert^2=2(1-\langle Qe,e\rangle),
		\qquad
		\langle Qe,u\rangle=1-\langle Qe,e\rangle,
	\]
	所以 \(r_u(Qe)=Qe-u=e\)。算子 \(r_uQ\) 固定 \(e\)，按上一段是至多 \(n-1\) 个反射的复合。由 \(Q=r_u(r_uQ)\)，总数至多为 \(n\)。

	对于仿射结论，先用\cref{g1:thm:isometry-normal}把变换写成 \(c+y\mapsto c+Qy+a\)，其中 \(a\in F=\ker(I-Q)\)。若 \(a=0\)，只须将线性结论中的反射超平面全部平移到经过 \(c\)，使用至多 \(n\) 个反射。若 \(a\ne0\)，令 \(e=a/\lVert a\rVert\)，在中心为 \(c\) 的坐标里，关于超平面
	\[
		H_s=\{y:\langle y,e\rangle=s\}
	\]
	的反射记为 \(\rho_s\)，其公式为 \(\rho_s(y)=y+2(s-\langle y,e\rangle)e\)。直接代入得 \(\rho_{\lVert a\rVert/2}\rho_0(y)=y+a\)，因此沿 \(a\) 的平移由两个平行超平面反射组成。另一方面，只须在 \(F^\perp\) 上分解 \(Q\)，再将反射扩充为在 \(F\) 上恒等，故需要至多 \(n-\dim F\) 个反射。因 \(a\ne0\) 保证 \(\dim F\geq1\)，总数不超过 \(n-\dim F+2\leq n+1\)。
\end{proof}

下面的低维分类来自两个已经证明的约束：正交块的维数相加必须等于空间维数，而不能消去的平移只能位于 \(+1\) 特征空间。保向与反向的情形都由这两个约束确定。以下“取坐标”均指更换原点和标准正交基；同一类中的角度与位移仍可能给出不同变换，其参数关系将在\secref{g1:sec:0104}中进一步讨论。

\ProseParagraphBreak

在一维，\(Q\) 只有 \(+1\) 与 \(-1\)。保向变换为 \(x\mapsto x+t\)，其中 \(t=0\) 是恒等，\(t\ne0\) 是非零平移。反向变换为 \(x\mapsto-x+t\)，有唯一不动点 \(t/2\)，即关于该点的反射。

\ProseParagraphBreak

在二维，保向正交算子或者为 \(I\)，或者是非恒等旋转。如果 \(Q=I\)，便得到恒等或非零平移；否则 \(F=0\)，所有平移都可消去，得到绕唯一中心的旋转，角度可以包含 \(\pi\)。反向正交算子只能有一个 \(+1\) 块与一个 \(-1\) 块，所以 \(F\) 是一条直线。不可消去的平移为零时得到关于一条直线的反射，非零时得到沿该直线的滑移反射。这五种互不重叠的情形列于\cref{g1:tab:plane-isometries}。

\begin{table}[htbp]
	\centering
	\caption{平面等距变换的类型}
	\label{g1:tab:plane-isometries}
	\begin{tabularx}{\linewidth}{@{}l l X@{}}
		\toprule
		类型 & 正规形及参数 & 不动点集 \\
		\midrule
		恒等 & \((x,y)\mapsto(x,y)\) & 整个平面 \\
		非零平移 & \((x,y)\mapsto(x+h,y),\ h>0\) & 空集 \\
		旋转 & \(v\mapsto R_\theta v,\ 0<\theta\leq\pi\) & 唯一中心 \\
		直线反射 & \((x,y)\mapsto(x,-y)\) & 反射轴 \\
		滑移反射 & \((x,y)\mapsto(x+h,-y),\ h>0\) & 空集 \\
		\bottomrule
	\end{tabularx}
\end{table}

二维表格允许翻转坐标轴方向，所以把非零平移长度取为正，把旋转角归入 \((0,\pi]\)。若坚持使用事先固定的平面定向，则旋转角应记作模 \(2\pi\) 的带方向角，不能仅凭该表消去正负。

\ProseParagraphBreak

在三维，保向而不恒等的正交算子必有形式 \(\operatorname{diag}(R_\theta,1)\)，其中允许 \(\theta=\pi\)。确实，若有旋转块，余下的一维块必须为 \(+1\)；若全部为 \(\pm1\) 块，行列式为正迫使 \(-1\) 块的数目为零或二，后一种也就是转角 \(\pi\) 的情形。因此保向等距变换只有恒等、非零平移、绕轴旋转和非平凡螺旋运动。后两种在 \(0<\theta\leq\pi\) 时有唯一的轴，即最小位移点集；\(h=0\) 与 \(h\ne0\) 将它们分开。

\ProseParagraphBreak

三维反向正交算子则可写成 \(\operatorname{diag}(R_\theta,-1)\)，其中 \(0\leq\theta\leq\pi\)。当 \(\theta=0\) 时，\(F\) 是二维平面，得到平面反射，或者平面反射后沿该平面的非零平移；后一种也称滑移反射。当 \(0<\theta\leq\pi\) 时，\(F=0\)，变换有唯一不动点，全部平移都能消去。此时它由绕一条轴的旋转和关于过中心、垂直于该轴的平面反射组成，称为\term[xuanzhuanfanshe]{旋转反射}[rotatory reflection]。这两步在坐标中分别作用于 \((x,y)\) 与 \(z\)，故可交换次序。当 \(\theta=\pi\) 时，正规形成为 \(-I\)，即\term[zhongxinduichen]{中心对称}[central inversion]；此时中心唯一，却没有唯一的旋转轴，任意过中心的直线都可用于上述分解。

\begin{table}[htbp]
	\centering
	\caption{三维等距变换的类型}
	\label{g1:tab:space-isometries}
	\begin{tabularx}{\linewidth}{@{}l l X@{}}
		\toprule
		类型 & 正规形及参数 & 不动点集 \\
		\midrule
		恒等 & \(v\mapsto v\) & 整个空间 \\
		非零平移 & \((x,y,z)\mapsto(x,y,z+h),\ h\ne0\) & 空集 \\
		绕轴旋转 & \((v,z)\mapsto(R_\theta v,z)\) & 旋转轴 \\
		螺旋运动 & \((v,z)\mapsto(R_\theta v,z+h),\ h\ne0\) & 空集 \\
		平面反射 & \((x,y,z)\mapsto(x,y,-z)\) & 反射平面 \\
		滑移反射 & \((x,y,z)\mapsto(x+h,y,-z),\ h>0\) & 空集 \\
		旋转反射 & \((v,z)\mapsto(R_\theta v,-z),\ 0<\theta<\pi\) & 唯一中心 \\
		中心对称 & \(v\mapsto-v\) & 唯一中心 \\
		\bottomrule
	\end{tabularx}
\end{table}

\cref{g1:tab:space-isometries}中的绕轴旋转与螺旋运动取 \(0<\theta\leq\pi\)，这两行的 \(v\) 是二维向量；恒等与中心对称两行的 \(v\) 是三维向量。表中将旋转反射的端点 \(0\) 和 \(\pi\) 分别归入平面反射与中心对称，将螺旋运动的 \(h=0\) 归入绕轴旋转，从而各行表示互不重叠的类型。纯平移的最小位移点集是整个空间，滑移反射的是反射平面，螺旋运动的是螺旋轴；这些集合比仅仅问“有没有不动点”提供了更细的几何信息。

% !TeX root = ../../Book1.tex
% 纲要标题：### 1.4 群作用如何组织几何分类
% 纲要位置：几何学计划纲要.md，第 144 行。
% 正文状态：初稿；尚未正式审核。

% 纲要内容开始
% 
% * 轨道、稳定子群与等价关系
% * 共轭、坐标选择和分类参数
% * Euclid 、仿射与相似分类的比较
% * 矩阵族退化时不变量的变化
% 
% ---
% 
% 纲要内容结束

\section{群作用如何组织几何分类}
\label{g1:sec:0104}

一个椭圆经刚性移动仍有原来的长短轴，经统一缩放会改变轴长，经一般仿射变换却可以成为圆。因此，分类之前必须先回答：哪些改变算作同一个对象？前面已经出现了仿射变换、正交变换和等距变换，现在把这些具体变换组织起来，就能说明何谓分类参数，以及参数为什么随允许的变换而改变。

\subsection{轨道、稳定子群与等价关系}
\label{g1:subsec:0104:01}

把可逆变换按复合相乘时，复合满足结合律，有恒等变换，每个变换也有逆变换。抽象地，一个集合及其乘法若满足结合律、存在单位元并且每个元素都有逆元，就称为\term[qun]{群}[group]。一个含单位元且对乘法、取逆封闭的子集称为\term[ziqun]{子群}[subgroup]。下面保留这些运算性质，而把变换实际作用的对象另记为 \(X\)。

\begin{definition}[群作用]\label{g1:def:group-action}
	设 \(G\) 为群，单位元为 \(e\)，\(X\) 为集合。 \(G\) 在 \(X\) 上的\term[qunzuoyong]{群作用}[group action]是运算 \((g,x)\mapsto g\cdot x\)，满足 \(e\cdot x=x\) 与 \((gh)\cdot x=g\cdot(h\cdot x)\)。定义
	\[
		G\cdot x\coloneqq\{g\cdot x:g\in G\},\qquad
		G_x\coloneqq\{g\in G:g\cdot x=x\},
	\]
	分别称为 \(x\) 的\term[guidao]{轨道}[orbit]与\term[wendingziqun]{稳定子群}[stabilizer subgroup]。
	\symidx{\(G\cdot x\)：点 \(x\) 的轨道}
	\symidx{\(G_x\)：点 \(x\) 的稳定子群}
\end{definition}

群在这里首先可以理解为一组可逆变换：乘法是复合，单位元是不改变对象的变换，逆元撤销原变换。稳定子群中的变换允许移动周围空间，但必须把指定对象留在原处；例如，一个球心固定的球仍容许许多旋转。

\begin{proposition}[轨道与稳定子群]\label{g1:prop:orbit-stabilizer}
	对任意群 \(G\) 在集合 \(X\) 上的作用，关系 \(x\sim y\Longleftrightarrow y\in G\cdot x\) 是等价关系，\(G_x\) 是 \(G\) 的子群。记 \(gG_x=\{gh:h\in G_x\}\)，称为子群 \(G_x\) 的\term[zuopeiji]{左陪集}[left coset]。左陪集集合 \(G/G_x=\{gG_x:g\in G\}\) 与轨道 \(G\cdot x\) 之间有双射
	\[
		gG_x\longmapsto g\cdot x.
	\]
	若 \(y=h\cdot x\)，则 \(G_y=hG_xh^{-1}\)。
\end{proposition}
\begin{proof}
	单位元给出自反性，逆元给出对称性，乘法给出传递性。若 \(g,h\) 都固定 \(x\)，则 \(gh\) 固定 \(x\)，而 \(g^{-1}\) 也固定 \(x\)，故 \(G_x\) 为子群。两个左陪集 \(gG_x,hG_x\) 相等，当且仅当 \(h^{-1}g\in G_x\)，这又等价于 \(g\cdot x=h\cdot x\)。所以所给映射既良定义又单射；满射由轨道定义得到。最后，\(k\) 固定 \(h\cdot x\) 当且仅当 \(h^{-1}kh\) 固定 \(x\)，这就是稳定子群的共轭公式。
\end{proof}

上述双射使用的是陪集集合 \(G/G_x\)；当 \(G_x\) 正规时，这个集合才带有自然的商群结构。若 \(G\) 有限，每个陪集的元素数等于 \(|G_x|\)，因而 \(|G\cdot x|=|G|/|G_x|\)。几何变换群通常无限，上面的集合双射仍然成立。

\begin{example}[点与有序标架]\label{g1:exm:frame-orbits}
	在 \(\mathbb R^n\) 中，所有可逆仿射映射组成\term[fangshequn]{仿射群}[affine group]，记为 \(\operatorname{Aff}(n,\mathbb R)\)。用 \((P,b)\) 表示 \(x\mapsto Px+b\)，其中 \(P\in\operatorname{GL}(n,\mathbb R)\)，则
	\[
		(P,b)(S,c)=(PS,Pc+b),\qquad
		(P,b)^{-1}=(P^{-1},-P^{-1}b).
	\]
	这些公式同时验证了复合与求逆仍在其中。这里 \(\operatorname{GL}(n,\mathbb R)\) 表示所有可逆 \(n\) 阶实矩阵组成的\term[yibanxianxingqun]{一般线性群}[general linear group]。仿射群在点集上的作用只有一个轨道，因为平移可以连接任意两点；原点的稳定子群为 \(\{(P,0)\}\)。它在有序仿射标架上的作用也只有一个轨道，但一个标架的稳定子群只含单位元：由\cref{g1:prop:frame-map}，标架的像唯一确定仿射映射。
\end{example}

分类需要\term[bubianliang]{不变量}[invariant]，即在同一轨道上取相同值的量。如果某个不变量还满足“取值相同就处于同一轨道”，便称它为\term[wanquanbubianliang]{完全不变量}[complete invariant]。秩是二次型合同变换的不变量，但实数域上仅知道秩不够： \(x_1^2+x_2^2\) 与 \(x_1^2-x_2^2\) 都有秩 \(2\)，惯性却不同。\cref{g1:thm:inertia}中的三个惯性数合在一起才给出完整分类。

\subsection{共轭、坐标选择和分类参数}
\label{g1:subsec:0104:02}

几何变换本身也可以作为待分类对象。若用可逆仿射映射 \(h(x)=Px+c\) 把旧坐标送到新坐标，计算新坐标中的作用时便先用 \(h^{-1}\) 还原旧坐标，再施行 \(f\)，最后用 \(h\) 返回新坐标。因此，\(f(x)=Qx+b\) 在新坐标中变成 \(hfh^{-1}\)。直接代入逆映射得
\begin{equation}\label{g1:eq:affine-conjugation}
	hfh^{-1}(y)=PQP^{-1}y+Pb+(I-PQP^{-1})c.
\end{equation}
这称为\term[gonge]{共轭}[conjugation]。取 \(P\) 为正交矩阵，便得到等距坐标改变下的公式。共轭允许换观察位置和方向，但保留迭代关系，因为 \((hfh^{-1})^k=hf^kh^{-1}\)。

\begin{proposition}[等距变换的共轭参数]\label{g1:prop:isometry-conjugacy}
	设 \(f_1,f_2\) 是同一 \(n\) 维实 Euclid 空间的等距变换。分别选取原点，将它们按\cref{g1:thm:isometry-normal}写成 \(f_i(x)=Q_ix+t_i\)，其中 \(t_i\in F_i\coloneqq\ker(I-Q_i)\)。则二者在 Euclid 等距变换群中共轭，当且仅当 \(Q_1,Q_2\) 正交共轭且 \(\lVert t_1\rVert=\lVert t_2\rVert\)。因此，正交标准分块中的 \(+1,-1\) 重数、\(0<\theta<\pi\) 的旋转角多重集，连同非负数 \(\lVert t_i\rVert\)，构成完全不变量。
\end{proposition}
\begin{proof}
	分别选择原点，等价于先把两个映射分别作一次平移共轭；如果变换后的映射可以相互共轭，前后复合这些平移及其逆就得到原映射间的共轭，反之亦然。因此可以比较各自的正规形。

	若 \(h(x)=Px+c\) 为共轭它们的等距变换，则 \(P\) 正交，且由\eqref{g1:eq:affine-conjugation}有 \(Q_2=PQ_1P^{-1}\)。所以 \(P(F_1)=F_2\)， \(Pt_1\in F_2\)。另一方面，\((I-Q_2)c\) 属于 \(F_2^\perp\)：对 \(v\in F_2\)， \(\langle(I-Q_2)c,v\rangle=\langle c,(I-Q_2^*)v\rangle=0\)。比较平移部分的 \(F_2\) 分量便得 \(t_2=Pt_1\)，从而长度相同。

	反过来，设正交映射 \(P_0\) 满足 \(P_0Q_1P_0^{-1}=Q_2\)。向量 \(P_0t_1\) 与 \(t_2\) 都在 \(F_2\) 且等长。若它们为零，取 \(U=I\)；否则把各自除以公共长度，再分别扩充为 \(F_2\) 的标准正交基，把前一组基映到后一组，得到 \(F_2\) 上的正交映射 \(U\)。令它在 \(F_2^\perp\) 上为恒等。由于 \(Q_2\) 在 \(F_2\) 上为恒等且保持其正交补， \(UQ_2=Q_2U\)。于是 \(P=UP_0\) 仍共轭两个正交部分，并满足 \(Pt_1=t_2\)。取 \(h(x)=Px\) 即得 \(hf_1h^{-1}=f_2\)。正交部分的参数由\cref{g1:thm:orthogonal-blocks}确定。
\end{proof}

这里允许所有等距坐标改变，含改变定向的变换。例如在一个二维旋转平面中翻转一个基向量，会把 \(R_\theta\) 变为 \(R_{-\theta}\)，所以参数可选在 \((0,\pi)\)。若只允许保定向的共轭，分类还可能需要额外的方向信息。

\ProseParagraphBreak

还要区分共轭与双线性型的坐标变换。同一个线性算子的输入与输出都换基，若旧坐标 \(x=Sx'\)，矩阵变成 \(S^{-1}AS\)；同一个二次型 \(x^{\mathsf T}Ax\) 改用新坐标则变成 \(x'^{\mathsf T}S^{\mathsf T}ASx'\)。前者是矩阵相似，后者是合同。矩阵相似是线性算子的换基；下一小节的几何相似变换则指统一缩放再作等距变换，名称相近而所讨论的对象不同。对正交矩阵，由 \(S^{-1}=S^{\mathsf T}\) 可知，两种公式在这个特殊条件下重合。

\subsection{Euclid 、仿射与相似分类的比较}
\label{g1:subsec:0104:03}

把正交条件加到\cref{g1:exm:frame-orbits}的仿射变换中，就得到\term[oushiqun]{Euclid 群}[Euclidean group] \(E(n)\)。它的元素为 \(x\mapsto Qx+b\)，由\cref{g1:thm:isometry-affine}也恰好是全部保持距离的自映射。进一步允许统一比例缩放，得到\term[xiangsiqun]{相似群}[similarity group] \(\operatorname{Sim}(n)\)，元素为 \(x\mapsto rQx+b\)，其中 \(r>0\)、\(Q\) 正交。这些映射把所有距离乘以同一个 \(r\)；反过来，具有这一距离性质的映射除去比例和平移后为等距映射。

\begin{table}[htbp]
	\centering\small
	\caption{三种变换所保留的信息}\label{g1:tab:transformation-groups}
	\begin{tabularx}{\linewidth}{@{}l l X@{}}
		\toprule
		变换 & 坐标形式 & 必定保留的几何量或关系 \\
		\midrule
		Euclid 等距变换 & \(Qx+b\) & 距离、角度、仿射组合与平行关系 \\
		相似变换 & \(rQx+b\) & 角度、非零长度之比、仿射组合与平行关系 \\
		可逆仿射变换 & \(Px+b\) & 仿射组合、子空间维数、相交与平行关系 \\
		\bottomrule
	\end{tabularx}
\end{table}

表中的包含关系为 \(E(n)\subseteq\operatorname{Sim}(n)\subseteq\operatorname{Aff}(n,\mathbb R)\)。允许的群越大，同一轨道可能包含越多对象，因而分类中能保留的信息越少。角度的保持可由内积公式核对：相似变换把内积与两个长度的乘积同时乘以 \(r^2\)，所以它们的比值不变。

\begin{example}[非退化椭球的三种分类]\label{g1:exm:ellipsoid-classification}
	在 \(\mathbb R^n\)、\(n\ge1\) 中，设 \(A\) 为实对称正定矩阵，考虑点集
	\[
		\mathcal E(A,c)=\{x:(x-c)^{\mathsf T}A(x-c)=1\}.
	\]
	\cref{g1:thm:spectral}将它化为 \(\sum_i\lambda_i y_i^2=1\)，其中 \(\lambda_i>0\)，半轴长为 \(a_i=\lambda_i^{-1/2}\)。在 Euclid 等距变换下，计入重数的无序半轴长多重集 \(\{a_1,\ldots,a_n\}\) 是完全不变量；在相似变换下，只剩它们的比例，即有序排列后归一化为最短半轴为 \(1\) 的长度组；在可逆仿射变换下，所有这样的椭球属于同一类。
\end{example}
\begin{proof}
	先证明椭球点集唯一确定中心和对称矩阵。正定性及谱分解说明该点集非空而有界。若 \(c,d\) 都是它的对称中心，两次中心对称的复合就是平移 \(2(d-c)\)。点集若在非零平移下不变，任取一点并反复平移便得到无界序列，与有界性矛盾。因此中心唯一。中心移到零后，每个单位向量 \(u\) 上的正射线与椭球相交于 \(u/\sqrt{u^{\mathsf T}Au}\)。所以点集确定 \(u^{\mathsf T}Au\)，由齐次性确定所有二次型值，再由极化公式确定对称矩阵 \(A\)。

	等距映射必须把唯一中心映到唯一中心，故中心化后两个矩阵通过正交合同联系，由谱定理的唯一性，特征值及半轴长的多重集一致。反过来，若半轴长组相同，匹配相应的标准正交特征基，再恢复中心的平移，就给出所求等距映射。相似变换则先把全部半轴长乘以公共正数 \(r\)，再作这种匹配，故判据恰为比例一致。最后，在正交特征坐标中置 \(z_i=\sqrt{\lambda_i}y_i\)，即可通过可逆线性变换把任意椭球送到单位球面。
\end{proof}

例如 \(x^2/4+y^2=1\) 与 \(x^2/16+y^2/4=1\) 相似而不等距，因为半轴由 \((2,1)\) 变为 \((4,2)\)。它们又都与圆仿射等价，但都不与圆相似，因为长短轴之比为 \(2\)。这也说明“能化成同一个标准方程”必须连同所用的坐标变换类别一起陈述。

\subsection{矩阵族退化时不变量的变化}
\label{g1:subsec:0104:04}

分类一个固定对象之后，还可以问：数据连续变化时，分类参数怎样变化？矩阵每个元素都很接近，并不保证它们处于同一个轨道。以下把矩阵族的连续性理解为各矩阵元的连续性；所用估计都在有限维空间内。

\begin{proposition}[秩与惯性的局部变化]\label{g1:prop:rank-inertia-continuity}
	设 \(A_0\) 为实矩阵。充分接近 \(A_0\) 的同型实矩阵 \(A\) 均满足 \(\operatorname{rank}A\ge\operatorname{rank}A_0\)。若矩阵皆为实对称 \(n\) 阶矩阵，记正、负惯性指数为 \(p(A),q(A)\)，则在 \(A_0\) 的某个邻域内
	\[
		p(A)\ge p(A_0),\qquad q(A)\ge q(A_0).
	\]
	所以在秩固定的连续实对称矩阵族中，惯性局部不变；特别地，非退化型的惯性局部不变。
\end{proposition}
\begin{proof}
	若 \(r=\operatorname{rank}A_0>0\)，取一个非零 \(r\) 阶子式。其行列式是矩阵元的多项式，故在充分小的邻域内仍非零，这证明秩的结论；\(r=0\) 时结论直接成立。

	由\cref{g1:thm:inertia}，有维数 \(p(A_0)\) 的子空间 \(U_+\)，使 \(u^{\mathsf T}A_0u\) 在 \(U_+\setminus\{0\}\) 上为正。若 \(p(A_0)>0\)，它在 \(U_+\) 的单位球面上的最小值 \(m_+\) 为正，因为该球面紧且型连续。对 \(D=A-A_0\)，用矩阵元的 Euclid 范数 \(\lVert D\rVert_{\mathrm F}=(\sum_{ij}d_{ij}^2)^{1/2}\) 有
	\[
		|u^{\mathsf T}Du|
		\le\left(\sum_{ij}d_{ij}^2\right)^{1/2}
		\left(\sum_{ij}u_i^2u_j^2\right)^{1/2}
		=\lVert D\rVert_{\mathrm F}\,\lVert u\rVert^2.
	\]
	所以当 \(\lVert D\rVert_{\mathrm F}<m_+/2\) 时，\(A\) 在同一个 \(U_+\) 上仍正定，从而 \(p(A)\ge p(A_0)\)；\(p(A_0)=0\) 时无需估计。在负定子空间上对 \(-A\) 作相同论证，得到负指数不等式。若同时有 \(\operatorname{rank}A=\operatorname{rank}A_0\)，利用 \(p(A)+q(A)=\operatorname{rank}A\)，两项不等式只能同时取等号。
\end{proof}

序列收敛时，极限处的秩可以降低；邻近矩阵的秩则至少等于极限矩阵的秩。这一性质称为秩的下半连续性。在一个参数区间上，秩固定就使惯性处处局部不变，因而在整个区间上不变：否则，取一个惯性值的原像以及其他所有惯性值原像的并；局部常值性说明二者都是开集，它们将该区间分成互不相交的两个非空开集，违背区间的连通性。

\begin{example}[越过退化型]\label{g1:exm:degenerating-form}
	实对称矩阵族 \(A_t=\operatorname{diag}(1,t)\) 对应 \(x^2+ty^2\)。当 \(t>0\) 时惯性为 \((2,0,0)\)，当 \(t=0\) 时为 \((1,0,1)\)，当 \(t<0\) 时为 \((1,1,0)\)。每个正参数都可用 \(z=\sqrt t\,y\) 化为 \(x^2+z^2\)，但这个变换在 \(t=0\) 时不可逆。因此，同一个非退化合同类的极限可以落入另一个轨道；分类所用的坐标不能保证跨过退化参数仍有效。
\end{example}

\begin{example}[特征值不变而相似类型改变]\label{g1:exm:shear-degeneration}
	考虑线性变换族
	\[
		N_t=\begin{pmatrix}1&t\\0&1\end{pmatrix}.
	\]
	当 \(t\ne0\) 时，取 \(P_t=\operatorname{diag}(t,1)\)，便有 \(N_t=P_tN_1P_t^{-1}\)，所以这些变换都相似。但 \(N_0=I\) 不与 \(N_1\) 相似：\(\operatorname{rank}(N_t-I)\) 在 \(t\ne0\) 时为 \(1\)，在 \(t=0\) 时为 \(0\)，而共轭保持该秩。全部特征值始终为 \(1\)，所以特征值多重集对一般线性算子并不是完全不变量。几何上，非零剪切只逐点固定一条直线，极限恒等变换却逐点固定整个平面。
\end{example}

谱定理中的自伴条件于是有了另一个可见作用：它排除了上述剪切型的现象，使特征值多重集足以完成正交分类。对于变化的几何对象，应分别观察连续数值参数、离散分类数据和达到退化点时失效的构造；不能仅凭某一个熟悉的不变量没有变化，就判断对象的类型没有变化。


\begin{chaptersummary}
	仿射空间把点与向量分开：点差是向量，系数和为 \(1\) 的组合才自然给出点。标架使这些运算可计算，但交集、方向空间与仿射包络不依赖坐标。加入内积后，正交补产生投影，自伴算子产生正交谱分解；一般对称型经合同变换只保留正、负、零惯性数，而 Gram 行列式同时给出体积的平方与线性依赖的判据。

	\ProseParagraphBreak

	距离保持足以推出 \(f(x)=Qx+b\)。正交部分分解为固定方向、反向方向与旋转平面；平移中垂直于固定空间的分量可以通过换原点消去，沿固定空间的分量则决定最小位移。低维的旋转、反射、滑移与螺旋运动是这一分解的具体形式。分类还须指定变换群：同一对象在 Euclid 、相似与仿射作用下有不同不变量；连续族达到退化点时，秩和惯性等离散信息可能改变，原先的标准化坐标也可能失效。
\end{chaptersummary}

\BookExercises

\begin{exercise}\label{g1:ex:prescribed-affine-map}
	在 \(\mathbb R^2\) 中，求将 \((0,0),(1,0),(0,1)\) 分别送到 \((1,1),(3,1),(2,4)\) 的仿射映射。求 \(p=(2,1)\) 的像和它相对于像标架的重心坐标。该映射是否为等距变换或相似变换？
\end{exercise}
\begin{solution}
	原点的像确定平移 \(b=(1,1)\)，两个基向量的像差确定矩阵的两列，故
	\[
		f(x,y)=\begin{pmatrix}2&1\\0&3\end{pmatrix}\binom{x}{y}+\binom11,
		\qquad f(2,1)=(6,4).
	\]
	其行列式为 \(6\ne0\)，所以像点组仍为标架。点 \(p\) 在原标架下的重心坐标为 \((-2,2,1)\)，因为这些坐标的和为 \(1\)，后两个系数是其仿射坐标。由仿射组合的保持性，同一重心坐标给出 \(f(p)\)，直接代入也得到 \(-2(1,1)+2(3,1)+(2,4)=(6,4)\)。原来两个垂直单位向量的像分别为 \((2,0),(1,3)\)，内积为 \(2\ne0\)，所以该映射既不等距也不相似。坐标与重心关系的保持并不要求角度保持。
\end{solution}

\begin{exercise}\label{g1:ex:parallelogram-midpoints}
	在实仿射空间中，设 \(p,q,r\) 仿射无关。证明恰有一个点 \(s\)，使 \(p,r\) 的中点等于 \(q,s\) 的中点，写出其仿射表达式，并说明 \(p,q,r,s\) 按此顺序是一个平行四边形。若不要求这三个点仿射无关，唯一性与平行四边形的非退化性中哪项仍成立？
\end{exercise}
\begin{solution}
	以 \(p\) 为原点，设 \(u=q-p,v=r-p,w=s-p\)。中点等式为 \(v/2=(u+w)/2\)，故唯一可能为 \(w=v-u\)，且它确实满足等式。所以 \(s=p-q+r\)，此式的系数和为 \(1\)，是合法仿射组合。对边向量分别满足 \(q-p=r-s=u\) 与 \(r-q=s-p=v-u\)。若 \(p,q,r\) 仿射无关，则 \(u,v\) 无关，\(u,v-u\) 也无关，四边形因此处于一个平面且不退化。无关性去掉后，方程仍唯一确定 \(s\)，但可能得到共线甚至重合的点，故非退化性不再成立。
\end{solution}

\begin{exercise}\label{g1:ex:affine-intersection-parameter}
	在 \(\mathbb R^3\) 中，令
	\[
		F=\{(s,0,0):s\in\mathbb R\},\qquad
		G_t=\{(0,u,t):u\in\mathbb R\}.
	\]
	分别求 \(F\cap G_t\) 与 \(\operatorname{aff}(F\cup G_t)\)，并说明只知道两条直线方向的夹角，为何不能判断它们是否相交。
\end{exercise}
\begin{solution}
	方程 \((s,0,0)=(0,u,t)\) 强制 \(s=u=t=0\)。所以 \(t=0\) 时交集为原点，\(t\ne0\) 时交集为空。包络的方向为 \(\operatorname{span}(e_1,e_2,te_3)\)，故 \(t=0\) 时包络为 \(z=0\) 平面， \(t\ne0\) 时为整个 \(\mathbb R^3\)。两条直线的方向始终垂直，发生变化的是基点差是否属于方向空间之和，夹角没有记录这一信息。
\end{solution}

\begin{exercise}\label{g1:ex:distance-configurations}
	设 \(p_0,\ldots,p_m\) 与 \(q_0,\ldots,q_m\) 是 \(\mathbb R^n\) 中两组有序点，满足所有两两距离 \(\lVert p_i-p_j\rVert=\lVert q_i-q_j\rVert\)。证明存在 Euclid 等距变换把每个 \(p_i\) 送到 \(q_i\)。若第一组点的仿射包络为 \(\mathbb R^n\)，证明此变换唯一。
\end{exercise}
\begin{solution}
	令 \(v_i=p_i-p_0,w_i=q_i-q_0\)。极化公式给出
	\[
		\langle v_i,v_j\rangle
		=\tfrac12\bigl(\lVert v_i\rVert^2+\lVert v_j\rVert^2-\lVert v_i-v_j\rVert^2\bigr)
		=\langle w_i,w_j\rangle.
	\]
	尝试定义 \(L(\sum_i a_iv_i)=\sum_i a_iw_i\)。若 \(\sum_i a_iv_i=0\)，由上述内积相等可知 \(\lVert\sum_i a_iw_i\rVert^2=\sum_{ij}a_ia_j\langle v_i,v_j\rangle=0\)，所以对应像也为零，定义因此与表示无关。同一计算表明 \(L\) 保内积，并把 \(\operatorname{span}(v_i)\) 双射到 \(\operatorname{span}(w_i)\)。取该子空间的标准正交基及其像，分别补成全空间的标准正交基，将补上的基向量逐一匹配，便将 \(L\) 扩充为正交算子 \(Q\)。所求映射为 \(x\mapsto q_0+Q(x-p_0)\)。若原点组仿射张成全空间，则这些 \(v_i\) 张成 \(\mathbb R^n\)，已唯一确定 \(Q\)，而 \(p_0\) 的像又唯一确定平移，故等距变换唯一。
\end{solution}

% 供章末习题整合；题目由本节概念独立设计。

\begin{exercise}
	\label{g1:ex:weighted-projection}
	在 \(\mathbb R^3\) 上取内积
	\[
		\langle x,y\rangle=x_1y_1+4x_2y_2+2x_3y_3.
	\]
	令 \(x=(2,0,1)\)。求 \(x\) 在直线 \(W=\mathbb R(1,1,0)\) 上的正交投影及到 \(W\) 的距离，再求它在仿射平面 \(F=\{\,y:y_1+y_2+y_3=1\,\}\) 上的最近点及到 \(F\) 的距离。说明为何不能直接把通常点积下的法向量 \((1,1,1)\) 用于第二问。
\end{exercise}

\begin{solution}
	令 \(v=(1,1,0)\)。由 \(\langle x,v\rangle=2\)、\(\lVert v\rVert^2=5\)，投影为
	\[
		P_Wx=\frac25v=(2/5,2/5,0).
	\]
	投影距离公式给出
	\[
		d(x,W)^2=\lVert x\rVert^2-\lVert P_Wx\rVert^2
		=6-\frac45=\frac{26}{5}.
	\]

	对平面，令 \(\alpha(y)=y_1+y_2+y_3\)。需要的是满足 \(\alpha(y)=\langle n,y\rangle\) 的向量，故 \(n=\alpha^\sharp=(1,1/4,1/2)\)，且 \(\lVert n\rVert^2=7/4\)。平面方程即 \(\langle n,y\rangle=1\)，而 \(\langle n,x\rangle=3\)。由\cref{g1:exm:plane-projection}的公式，最近点为
	\[
		x-\frac{3-1}{7/4}n
		=\left(\frac67,-\frac27,\frac37\right),
	\]
	其坐标和等于 \(1\)。距离平方为
	\[
		\left(\frac87\right)^2\frac74=\frac{16}{7}.
	\]
	故两问距离分别为 \(\sqrt{26/5}\) 与 \(4/\sqrt7\)。向量 \((1,1,1)\) 在指定内积下对应的线性函数是 \(y_1+4y_2+2y_3\)，并非平面的定义函数；它甚至不与平面方向向量 \((1,-1,0)\) 正交，二者内积为 \(-3\)。
\end{solution}

\begin{exercise}
	\label{g1:ex:simultaneous-spectral}
	设 \(S,T\) 是同一个实有限维内积空间上的自伴算子。证明：存在同时由 \(S\)、\(T\) 的特征向量组成的标准正交基，当且仅当 \(ST=TS\)。再用
	\[
		S=\begin{pmatrix}1&0\\0&-1\end{pmatrix},\qquad
		T=\begin{pmatrix}0&1\\1&0\end{pmatrix}
	\]
	说明两个算子各自可以正交对角化，并不保证它们可以在同一组基中对角化。
\end{exercise}

\begin{solution}
	若有共同特征基，则两个矩阵在这组基下都为对角矩阵，故可交换。

	反过来，设 \(ST=TS\)。由\cref{g1:thm:spectral}，空间是 \(S\) 的各特征子空间 \(E_\lambda(S)\) 的正交直和。若 \(v\in E_\lambda(S)\)，则
	\[
		S(Tv)=T(Sv)=\lambda Tv,
	\]
	所以 \(Tv\in E_\lambda(S)\)。因此可以将 \(T\) 限制到每个 \(E_\lambda(S)\)，限制仍自伴：两边内积只是原内积的限制，自伴等式不变。再次应用谱定理，在每个 \(E_\lambda(S)\) 中取 \(T\) 的标准正交特征基。空间之间互相正交，合并这些基得到整个空间的标准正交基，其中每个向量又都是 \(S\) 的特征向量。

	题给矩阵均为实对称矩阵，但
	\[
		ST=\begin{pmatrix}0&1\\-1&0\end{pmatrix},\qquad
		TS=\begin{pmatrix}0&-1\\1&0\end{pmatrix}.
	\]
	两者不同，故不存在共同标准正交特征基。
\end{solution}

\begin{exercise}
	\label{g1:ex:inertia-family}
	对实参数 \(t\)，求二次型
	\[
		q_t(x,y,z)=x^2+2txy+y^2+z^2
	\]
	的惯性、根基与非退化参数范围。给出一组对所有 \(t\) 都适用的正交坐标，并说明惯性在哪些参数处发生变化。
\end{exercise}

\begin{solution}
	取 \(u=(x+y)/\sqrt2\)、\(v=(x-y)/\sqrt2\)、\(w=z\)，这是与 \(t\) 无关的正交换基，给出
	\[
		q_t=(1+t)u^2+(1-t)v^2+w^2.
	\]
	当 \(-1<t<1\) 时三个系数都正，惯性为 \((3,0,0)\)。当 \(t>1\) 或 \(t<-1\) 时恰有一个系数为负，惯性为 \((2,1,0)\)。当 \(t=1\) 或 \(t=-1\) 时恰有一个系数为零，其余两个为正，惯性为 \((2,0,1)\)。

	在 \(t\ne\pm1\) 时根基为 \(\{0\}\)。在 \(t=1\) 时根基是原坐标中的 \(\mathbb R(1,-1,0)\)；在 \(t=-1\) 时根基是 \(\mathbb R(1,1,0)\)。因此非退化参数恰为 \(\mathbb R\setminus\{-1,1\}\)。惯性在这两个参数处改变，对应的是一个特征值经过零；在三个互不相连的非退化参数区间内，它保持不变。
\end{solution}

\begin{exercise}
	\label{g1:ex:equidistant-four-points}
	设 Euclid 空间中四点 \(p_0,p_1,p_2,p_3\) 两两距离都是 \(1\)。不用预先画出正四面体，证明它们不能同处于一个仿射平面，并在 \(\mathbb R^3\) 中给出一组满足条件的坐标。提示：令 \(v_i=p_i-p_0\)，由距离恢复 \(\langle v_i,v_j\rangle\)。
\end{exercise}

\begin{solution}
	由 \(\lVert v_i\rVert=1\)，且 \(\lVert v_i-v_j\rVert=1\)，展开平方范数得
	\[
		\langle v_i,v_j\rangle
		=\frac{\lVert v_i\rVert^2+\lVert v_j\rVert^2-\lVert v_i-v_j\rVert^2}{2}
		=\frac12\qquad(i\ne j).
	\]
	故三条向量的 Gram 矩阵为
	\[
		\Gamma=\begin{pmatrix}1&1/2&1/2\\1/2&1&1/2\\1/2&1/2&1\end{pmatrix},
		\qquad \det\Gamma=\frac12>0.
	\]
	由\cref{g1:prop:gram-volume}，它们线性无关，四点的仿射张成空间因而是三维的，不能包含于一个仿射平面。

	在 \(\mathbb R^3\) 中可以取
	\[
		p_0=(0,0,0),\quad p_1=(1,0,0),\quad
		p_2=(1/2,\sqrt3/2,0),\quad
		p_3=(1/2,1/(2\sqrt3),\sqrt{2/3}).
	\]
	后三个位置向量的长度平方分别为 \(1\)，两两内积都是 \(1/2\)：其中 \(\langle p_2,p_3\rangle=1/4+1/4=1/2\)。于是它们与 \(p_0\) 的距离为 \(1\)，其余两两距离平方为 \(1+1-2(1/2)=1\)。这验证了全部六条距离条件。
\end{solution}

\begin{exercise}
	\label{g1:ex:positive-square-root}
	设 \(T\) 是实有限维内积空间 \(V\) 上的自伴算子，且 \(\langle Tv,v\rangle>0\) 对所有 \(v\ne0\) 成立。证明存在唯一的自伴正定算子 \(R\) 满足 \(R^2=T\)。说明若删除对 \(R\) 的正定要求，唯一性为何失败；再说明内积 \(B(u,v)=\langle Tu,v\rangle\) 与原内积之间有什么关系。
\end{exercise}

\begin{solution}
	由谱定理，\(V\) 是 \(T\) 的各特征子空间 \(E_\lambda(T)\) 的正交直和，且全部特征值 \(\lambda>0\)。在 \(E_\lambda(T)\) 上定义 \(R=\sqrt\lambda I\)，再按直和延拓。这个算子在标准正交特征基下是正对角矩阵，所以自伴、正定，且逐个特征子空间上都有 \(R^2=T\)。

	若 \(S\) 是另一个自伴正定平方根，则 \(ST=SS^2=S^3=S^2S=TS\)，故 \(S\) 保持每个 \(E_\lambda(T)\)。限制到这个子空间，\(S\) 仍自伴、正定，因而可由谱定理正交对角化。它的每个特征值 \(\mu\) 都满足 \(\mu>0\) 和 \(\mu^2=\lambda\)，所以只能等于 \(\sqrt\lambda\)。于是 \(S\) 在该子空间上等于 \(\sqrt\lambda I\)，即等于 \(R\)。各特征子空间张成 \(V\)，故 \(S=R\)。

	若 \(V\ne\{0\}\)，删除正定条件后，\(R\) 和 \(-R\) 是两个不同的自伴平方根；还可以在不同特征方向上选择不同正负号。最后，由 \(R\) 自伴及 \(R^2=T\)，
	\[
		B(u,v)=\langle R^2u,v\rangle=\langle Ru,Rv\rangle.
	\]
	因此可逆线性映射 \(R\) 把带内积 \(B\) 的 \(V\) 等距地映到带原内积的 \(V\)。它解释了为什么任意正定型都能经线性坐标伸缩变成平方和；该伸缩相对于原内积一般不保持长度。
\end{solution}


% 供第一章习题整合；以下均不承担正文证明责任。
\begin{exercise}
	\label{g1:ex:isometry-frame}
	设 \(f\colon\mathbb R^2\rightarrow\mathbb R^2\) 保持距离，并且
	\[
		f(0,0)=(1,2),\qquad f(1,0)=(1,3),\qquad f(0,1)=(0,2).
	\]
	求 \(f\) 的公式，判断它的类型，并求不动点。
\end{exercise}
\begin{solution}
	由\cref{g1:thm:isometry-affine}，\(f(x)=Qx+t\)，其中 \(t=(1,2)\)。减去 \(f(0,0)\) 后，两列分别为 \(Q(1,0)=(0,1)\)、\(Q(0,1)=(-1,0)\)，故
	\[
		f(x,y)=(-y+1,x+2).
	\]
	它的正交部分为旋转 \(\pi/2\)，固定方向空间为零，因而有唯一不动点。直接解 \(x=-y+1\)、\(y=x+2\)，得到 \((-1/2,3/2)\)。以此为原点便是逆时针 \(\pi/2\) 旋转。这里三个点的像决定公式，是因为其两个差向量构成一组基，不能只凭“给了三个条件”来判断唯一性。
\end{solution}

\begin{exercise}
	\label{g1:ex:space-glide}
	设 \(f\colon\mathbb R^3\rightarrow\mathbb R^3\) 为
	\[
		f(x,y,z)=(y+1,x-1,z+2).
	\]
	求其最小位移及取得最小位移的点集，说明它的几何类型，并求 \(f^2\)。
\end{exercise}
\begin{solution}
	正交部分为 \(Q(x,y,z)=(y,x,z)\)，所以
	\[
		F=\{(x,y,z):x=y\},\qquad F^\perp=\mathbb R(1,-1,0).
	\]
	平移向量分解为 \(t=(0,0,2)+(1,-1,0)\)，因此 \(a=(0,0,2)\)。取 \(c=(1/2,-1/2,0)\)，有 \((I-Q)c=(1,-1,0)\)。由\cref{g1:thm:isometry-normal}，最小位移为 \(2\)，在平面 \(c+F=\{x-y=1\}\) 上且仅在其上取得。\(Q\) 在 \(F\) 上恒等、在其法向上反号，故 \(f\) 是关于平面 \(x-y=1\) 的反射，再沿该平面平移 \((0,0,2)\)，属于反向的滑移反射。螺旋运动的非恒等正交部分保向，并且最小位移点集是一条直线，两者在此已可区分。最后直接计算
	\[
		f^2(x,y,z)=f(y+1,x-1,z+2)=(x,y,z+4).
	\]
\end{solution}

\begin{exercise}
	\label{g1:ex:two-reflections}
	在已固定标准定向的 Euclid 平面中，记 \(r_\alpha\) 为关于过原点且方向角为 \(\alpha\) 的直线的反射。
	\begin{enumerate}
		\item 求 \(r_\beta r_\alpha\)，并说明复合次序对结果的影响。
		\item 记 \(s_a\) 为关于直线 \(x=a\) 的反射。求 \(s_b s_a\)，并求 \(a\ne b\) 时的位移长度。
	\end{enumerate}
\end{exercise}
\begin{solution}
	1. 关于 \(x\) 轴的反射矩阵为 \(D=\operatorname{diag}(1,-1)\)。先旋转 \(-\alpha\) 使反射轴成为 \(x\) 轴，再反射，最后旋转回去，得到
	\[
		r_\alpha=R_\alpha D R_{-\alpha}
		=\begin{pmatrix}
			\cos2\alpha&\sin2\alpha\\
			\sin2\alpha&-\cos2\alpha
		\end{pmatrix}.
	\]
	由 \(DR_\gamma=R_{-\gamma}D\)，
	\[
		r_\beta r_\alpha
		=R_\beta D R_{\alpha-\beta}D R_{-\alpha}
		=R_{2(\beta-\alpha)}.
	\]
	这里右侧映射先作用，所以先关于方向角 \(\alpha\) 的直线反射，再关于方向角 \(\beta\) 的直线反射。交换次序会把转角改为 \(2(\alpha-\beta)\)，即得到逆旋转。反射轴的方向角只在模 \(\pi\) 意义下确定，但双倍角在模 \(2\pi\) 意义下不受此选择影响。

	2. 由于 \(s_a(x,y)=(2a-x,y)\)，有
	\[
		s_b s_a(x,y)=(2b-2a+x,y).
	\]
	因此复合为沿两条直线的公共法向平移 \(2(b-a)\)，位移长度为 \(2|b-a|\)，等于两条平行线距离的两倍。
\end{solution}

\begin{exercise}
	\label{g1:ex:minimal-reflections}
	设 \(Q\) 是 \(n\) 维实内积空间的正交算子。证明：把 \(Q\) 表示为过原点的超平面反射的复合，所需反射的最少个数等于 \(\operatorname{rank}(I-Q)\)。
\end{exercise}
\begin{solution}
	先证下界。任意同维线性算子 \(A,B\) 满足
	\[
		I-AB=(I-A)+A(I-B),
	\]
	故 \(\operatorname{rank}(I-AB)\leq\operatorname{rank}(I-A)+\operatorname{rank}(I-B)\)。若 \(r\) 为一个超平面反射，则 \(I-r\) 的像恰为法向直线，秩为 \(1\)。反复使用上述不等式，若 \(Q=r_1\cdots r_m\)，便得到 \(\operatorname{rank}(I-Q)\leq m\)。

	再证该下界能够达到。令 \(F=\ker(I-Q)\)，并设 \(d=\dim F^\perp\)。秩与核维数公式给出 \(d=\operatorname{rank}(I-Q)\)。空间 \(F^\perp\) 被 \(Q\) 保持，限制在其上的算子仍正交；由\cref{g1:thm:reflection-generation}，它是至多 \(d\) 个 \(F^\perp\) 内的超平面反射的复合。将每个反射按同一法向量扩充到整个空间，便在 \(F\) 上恒等，复合得到 \(Q\)。所需个数至多 \(d\)，与下界合起来恰为 \(d\)。\(d=0\) 时 \(Q=I\)，与零次反射的约定一致。
\end{solution}

\begin{exercise}
	\label{g1:ex:isometry-bounded-orbit}
	设 \(f\) 是 \(\mathbb R^n\) 的等距变换。对一点 \(x\)，其正向轨道指集合 \(\{f^k(x):k=0,1,2,\ldots\}\)。
	\begin{enumerate}
		\item 证明：\(f\) 有一条有界正向轨道，当且仅当 \(f\) 有不动点；在此情形，每条轨道都有界。
		\item 证明：存在正整数 \(m\) 使 \(f^m=\operatorname{id}\)，当且仅当 \(f\) 的不可消去平移分量为零，且正交部分各旋转块的角度除以 \(2\pi\) 都是有理数。
		\item 给出有界轨道不蕴含变换有限阶的具体例子。
	\end{enumerate}
\end{exercise}
\begin{solution}
	1. 按\cref{g1:thm:isometry-normal}，写 \(x=c+u+v\)，其中 \(u\in F\)、\(v\in F^\perp\)，则归纳得到
	\[
		f^k(x)=c+u+ka+Q^kv.
	\]
	其在 \(F\) 上相对于 \(c\) 的投影为 \(u+ka\)。若 \(a\ne0\)，该投影长度随 \(k\) 无界，所以轨道不可能有界。若 \(a=0\)，则 \(\lVert f^k(x)-(c+u)\rVert=\lVert v\rVert\)，每条轨道都有界；同一定理又说明 \(a=0\) 当且仅当 \(f\) 有不动点。

	2. 若 \(f^m=\operatorname{id}\)，则每条轨道都是有限集，由第1问知 \(a=0\)。以不动点为原点后 \(f=Q\)，且 \(Q^m=I\)。对一个旋转块，\(R_\theta^m=R_{m\theta}\)，故 \(m\theta\) 是 \(2\pi\) 的整数倍，\(\theta/(2\pi)\) 为有理数。反之，若 \(a=0\) 且全部旋转角满足此条件，取这些有理数分母的公倍数，再乘以 \(2\)，得到正整数 \(m\)。每个旋转块的 \(m\) 次方为 \(I\)，\(-1\) 块的 \(m\) 次方也为 \(1\)，而 \(+1\) 块始终不变，因此 \(Q^m=I\)，进而 \(f^m=\operatorname{id}\)。没有旋转块时直接取 \(m=2\) 即可。

	3. 平面上绕原点旋转角 \(2\pi\sqrt2\) 的变换，保持每个点到原点的距离，所以轨道都有界。若某个正整数 \(m\) 使其 \(m\) 次方为恒等，就会有 \(m\sqrt2\in\mathbb Z\)，与 \(\sqrt2\) 无理矛盾。因此它不是有限阶变换。
\end{solution}


\begin{exercise}\label{g1:ex:ellipsoid-stabilizer}
	设 \(A=\operatorname{diag}(\lambda_1,\ldots,\lambda_n)\) 为正定实对称矩阵。证明保持椭球 \(x^{\mathsf T}Ax=1\) 的 Euclid 等距变换恰为满足 \(QA=AQ\) 的正交变换。若不同特征值的重数为 \(m_1,\ldots,m_k\)，描述其稳定子群；比较所有半轴长互异与所有半轴长相等的情形。
\end{exercise}
\begin{solution}
	由\cref{g1:exm:ellipsoid-classification}的中心唯一性，保持点集的等距变换必须固定原点，所以为 \(x\mapsto Qx\)。该证明还说明中心化椭球的点集唯一决定二次型，故保持点集等价于 \(Q^{\mathsf T}AQ=A\)。由于 \(Q^{-1}=Q^{\mathsf T}\)，这等价于 \(AQ=QA\)。交换关系说明 \(Q\) 保持每个特征空间，反之在各特征空间上独立选择正交变换便与 \(A\) 交换。因此稳定子群同构于 \(O(m_1)\times\cdots\times O(m_k)\)。半轴互异时每个空间一维，只能独立翻转各轴，群共有 \(2^n\) 个元素；全部半轴相等时椭球为球面，稳定子群为整个 \(O(n)\)。
\end{solution}

\begin{exercise}\label{g1:ex:affine-versus-euclidean}
	在 \(\mathbb R^2\) 中考虑非零平移 \(\tau_v(x)=x+v\)。证明任意两个非零平移在仿射群中共轭，而在 Euclid 群中共轭当且仅当平移向量等长。再说明面积保持并不足以推出一个线性变换是等距变换。
\end{exercise}
\begin{solution}
	代入\eqref{g1:eq:affine-conjugation}中 \(Q=I,b=v\)，得到 \(h\tau_vh^{-1}=\tau_{Pv}\)。对任意非零 \(v,w\)，分别将它们扩充成一组基，将两组基匹配便得到可逆 \(P\) 满足 \(Pv=w\)，故仿射共轭成立。当 \(P\) 必须正交时，必要条件为 \(\lVert v\rVert=\lVert w\rVert\)；等长时匹配以归一化向量为首向量的标准正交基，又给出充分性。最后，剪切 \((x,y)\mapsto(x+y,y)\) 行列式为 \(1\)，保持平行四边形面积，却把单位向量 \((0,1)\) 变为长度 \(\sqrt2\) 的向量，故不等距。
\end{solution}

\begin{exercise}\label{g1:ex:orbit-closure}
	令 \(D=\operatorname{diag}(1,-1)\)。在实对称 \(2\) 阶矩阵空间中，构造 \(S_t\in\operatorname{GL}(2,\mathbb R)\)、\(t>0\)，使 \(S_t^{\mathsf T}DS_t\) 当 \(t\rightarrow0\) 时分别趋于 \(\operatorname{diag}(1,0)\)、\(\operatorname{diag}(0,-1)\) 与零矩阵。证明这同一个合同轨道中的矩阵序列不可能趋于单位矩阵。
\end{exercise}
\begin{solution}
	分别取 \(S_t=\operatorname{diag}(1,t)\)、\(\operatorname{diag}(t,1)\) 与 \(tI\)，所得矩阵为 \(\operatorname{diag}(1,-t^2)\)、\(\operatorname{diag}(t^2,-1)\) 与 \(t^2D\)，极限如题。三个极限的惯性都与 \(D\) 不同，故轨道不包含它们。若存在轨道内矩阵 \(B_j\rightarrow I\)，则当 \(j\) 充分大时 \(\lVert B_j-I\rVert_{\mathrm F}<1/2\)，从而对非零 \(x\) 有 \(x^{\mathsf T}B_jx\ge(1-1/2)\lVert x\rVert^2>0\)。这与轨道中每个矩阵的负惯性指数为 \(1\) 矛盾。因此轨道闭包增加的是某些退化型，不能任意增加别的非退化惯性类型。
\end{solution}
